\documentclass[journal]{IEEEtran}

\usepackage[T1]{fontenc}
\usepackage{amsmath,amssymb,amsfonts}
\usepackage{mathtools}
\usepackage{bm}
\usepackage{booktabs}
\usepackage{multirow}
\usepackage{array}
\usepackage{tabularx}
\usepackage{graphicx}
\usepackage{cite}
\usepackage{url}
\usepackage{float}
\usepackage{stfloats}
\usepackage{placeins}
\usepackage{xcolor}
\usepackage[hidelinks]{hyperref}
\usepackage{tikz}
\usetikzlibrary{calc, shapes, arrows.meta, positioning, backgrounds}
\usepackage{siunitx}
\AtBeginDocument{
    \RenewCommandCopy\qty\SI
}
% --- Define \subparagraph for IEEEtran ---
\makeatletter
\newcommand{\subparagraph}{\@startsection{subparagraph}{5}{\parindent}%
  {0.25ex plus 0.1ex minus 0.1ex}%  % Vertical space above heading
  {-0.5em}%                         % Negative value creates a run-in heading (horizontal space after)
  {\normalfont\normalsize\itshape}} % Heading typography (italicized, normal size)
\makeatother

\hypersetup{
    colorlinks=true,
    linkcolor=blue,
    citecolor=blue,
    urlcolor=blue
}

\title{
JANUS: A Constraint-Aware Hybrid Photonic Architecture with One-Hot Optical RNS Multiplication, High-Power Passive Fan-Out, and Bounded Exact-Precision Computation
}

\author{
\IEEEauthorblockN{Deepanshu Bhardwaj}\\[2pt]
\IEEEauthorblockA{\small Patent Application No.: 202611052791 (Patent Pending)}
}

\begin{document}
\maketitle

\begin{abstract}

Large-scale photonic accelerators are fundamentally constrained by analog precision collapse, optical loss accumulation, thermal drift, nonlinear wavelength interference, and the unsustainable static power requirements of thermo-optic routing meshes. Conventional photonic matrix multiplication architectures rely upon analog amplitude accumulation followed by high-speed analog-to-digital conversion, implicitly assuming that signal fidelity can scale with matrix dimension. However, for a $128 \times 128$ photonic multiply-accumulate operation using 8-bit operands, the accumulated output spans more than $8.3 \times 10^{6}$ discrete levels, requiring approximately $138$~dB signal-to-noise ratio (SNR) for lossless analog recovery---a precision-throughput tradeoff that becomes increasingly impractical under realistic shot noise, thermal noise, and ADC bandwidth constraints at multi-GHz operating frequencies. Mitigation techniques exist, but each substantially reduces effective throughput or efficiency. 

JANUS resolves this contradiction by abandoning analog amplitude reconstruction entirely. Instead of encoding numerical magnitude in optical intensity, JANUS reformulates computation as a spatial localization problem using One-Hot Optical Residue Number System (RNS) routing \cite{Omosanya2021}\cite{Huang1979}\cite{Sorger2018}. Each computation is represented as the position of light within spatial waveguides, reducing photodetection to a binary decision: whether light exists within a single channel. This substantially reduces the required detection precision from a $138$~dB analog precision target toward a binary threshold detection problem: whether signal power exceeds noise power in a single spatial channel. Binary threshold detection requires substantially lower detector precision and dynamic range than high-resolution analog reconstruction, while preserving exact integer residue computation.

To compensate for passive fan-out losses without introducing amplified spontaneous emission (ASE) noise from intermediate optical gain elements, the architecture integrates a high-power \SI{1064}{\nano\meter} ytterbium-fiber laser source paired with high-speed Separate-Absorption-Charge-Cliff-Multiplication ($\mathrm{SAC^2M}$) Ge/Si avalanche photodetectors (APDs) \cite{Jeong2004}\cite{Chen2025}. A non-volatile Asymmetric 16-Tree Fermat routing fabric constructed from Phase-Change Material (PCM) switches ($\mathrm{Sb_2S_3}$) with the 16th waveguide ($WG_{16}$) extension achieves an ultra-shallow depth of only 4 binary stages and 240 switches per multiplier \cite{Wuttig2017}\cite{Delaney2020}\cite{Rios2021}, restricting fabric insertion loss to only $1.61~\mathrm{dB}$ while eliminating static hold power. The resulting 17-channel spatial alphabet ($WG_0 \dots WG_{16}$) natively realizes the Fermat Prime field $\mathbb{Z}_{17}$ with $100\%$ state efficiency ($289/289$ states) and unlocks Radix-16 $\mathbb{Z}_{257}$ division-free modular reduction via Tree 16 negation symmetry ($16 \times W \equiv -W \pmod{17}$). Single-wavelength coherent fanout substantially suppresses dominant inter-channel nonlinear interference mechanisms, including four-wave mixing (FWM) and cross-phase modulation (XPM), by architectural construction rather than post-hoc filtering.

Each optical tile natively implements a modular residue channel over spatial waveguides, natively executing Fermat Prime arithmetic ($\mathbb{Z}_{17}$ and $\mathbb{Z}_{257}$) with division-free reduction. To scale precision to wide integer word lengths (e.g., INT8 to INT64), JANUS employs Spatial Residue Number System decomposition across $k$ parallel tiles and exact digital reconstruction using the Chinese Remainder Theorem (CRT) \cite{Mohan2016}\cite{Wang2000}\cite{Xia2000}. This decouples arithmetic precision from analog optical fidelity, operating entirely within bounded thermal, optical, and electrical regimes while sustaining deterministic exact computation. 

Across three deployment configurations---Mini, Edge, and Datacenter---JANUS achieves sustained exact 64-bit throughput ranging from $174.1$~TMAC/s to $2,785.3$~TMAC/s under bounded operating conditions without continuous optical calibration. The architecture demonstrates that large-scale photonic AI acceleration becomes physically realizable only when analog precision recovery is replaced by spatially encoded exact arithmetic.

\end{abstract}

\begin{IEEEkeywords}
Photonic Computing, Optical Neural Networks, Residue Number Systems, One-Hot Encoding, Phase-Change Materials, Avalanche Photodetection, Passive Fan-Out, Ytterbium- Doped Fiber Lasers, Sb2S3, Silicon Photonics, Thermal Management
\end{IEEEkeywords}

\section{INTRODUCTION}

The rapid scaling of artificial intelligence workloads has intensified interest in photonic computing architectures capable of exceeding the bandwidth-density and energy-efficiency limits of conventional electronic accelerators \cite{Miller2017}\cite{Miller2013}. Optical systems naturally provide high-degree spatial parallelism, low propagation delay, and wavelength-domain multiplexing, making silicon photonics an attractive substrate for large-scale matrix multiplication and neural network acceleration \cite{Shen2017}\cite{Tait2017}\cite{Ashtiani2022}\cite{Feldmann2019}. 

Most contemporary photonic neural accelerators are constructed around Mach-Zehnder Interferometer (MZI) meshes performing analog optical multiply-accumulate operations \cite{Shen2017}\cite{Harris2018}\cite{Feldmann2021}\cite{Reed2010}\cite{Sinatkas2021}. In these systems, numerical values are encoded into optical amplitudes, propagated through programmable interferometric meshes, and accumulated as continuous analog optical power before photodetection and analog-to-digital conversion.

While this paradigm demonstrates promising low-bit inference capability, it encounters severe scaling contradictions when extended toward large-scale exact computation. The contradiction originates from a fundamental mismatch between analog optical signal fidelity and the information density required by high-dimensional arithmetic accumulation.

Consider a conventional optical matrix multiplication involving 128 independent 8-bit activations and 128 independent 8-bit weights. The maximum unreduced analog accumulation is

\begin{align}
S_{\max} &= 128 \times 255 \times 255 \\
&= 8{,}323{,}200 \label{eq:max_accumulation}
\end{align}

Exact recovery of this value through analog photodetection requires resolution across more than $8.3 \times 10^{6}$ discrete analog levels. Under standard quantization theory, the minimum required SNR becomes \cite{Winzer2006,Saleh2019}

\begin{align}
\mathrm{SNR}_{req} &= 20\log_{10}\left( 8{,}323{,}200 \right) \\
&\approx 138.4~\mathrm{dB} \label{eq:snr_requirement}
\end{align}

This requirement exceeds the practical operating range of modern photonic systems by a substantial margin. Shot noise, thermal noise, amplifier noise, detector nonlinearity, and finite ADC precision collectively prevent stable analog recovery at such information densities.
The failure is not merely an implementation limitation.

It is a structural consequence of attempting to scale analog amplitude accumulation beyond physically sustainable entropy density.

Several mitigation strategies have been explored within prior work, including signal normalization, temporal bit slicing, calibration loops, and multi-stage electronic correction \cite{Hamerly2022}, \cite{Cauwenberghs2013}. However, these approaches only redistribute the precision
burden rather than eliminating it. Temporal decomposition reduces throughput proportionally to bit depth, while aggressive normalization sacrifices dynamic range without solving the underlying Shannon-capacity limitation. Consequently, existing photonic accelerators remain constrained to low-bit approximate inference regimes.

JANUS originates from the observation that the analog accumulation model itself constitutes the primary scaling failure. Rather than attempting to improve analog precision recovery, JANUS reformulates optical computation into a spatial localization problem.

The central architectural transition is therefore: 

\textbf{Amplitude Reconstruction $\Longrightarrow$ Spatial Localization}

Instead of representing values through optical intensity, JANUS represents values through spatial waveguide position using One-Hot Optical Residue Number System (RNS) encoding. Each input activates exactly one waveguide among the spatial channels ($N \le 256$ in general RNS formulations, or 17 channels in the Asymmetric 16-Tree Fermat Core). Photodetection no longer estimates analog amplitude. Instead, the detector answers a binary question:

\begin{center}
\textit{``Does waveguide k contain light?''}
\end{center}

This transforms the photodetection problem from high resolution analog conversion into binary spatial detection. Binary threshold detection requires substantially lower detector precision and dynamic range than high-resolution analog reconstruction. The required detector condition therefore reduces to verifying that signal power exceeds the noise floor within the active channel:

\begin{align}
\mathrm{SNR}_{\mathrm{binary}}
&>
\mathrm{Noise\ Floor}
\label{eq:binary_requirement}
\end{align}

By substantially reducing the analog precision requirement, JANUS suppresses the dominant precision bottleneck preventing large-scale photonic exact arithmetic.

However, resolving the precision contradiction alone does not produce a physically realizable architecture. Large-scale photonic systems simultaneously encounter several additional scaling failures:

\begin{itemize}
\item  Thermo-optic MZI routing meshes consume static hold power proportional to switch count, producing unsustainable thermal densities at scale.
\item  Dense wavelength-division multiplexing (DWDM) architectures suffer from nonlinear interference mechanisms including four-wave mixing (FWM) and cross-phase modulation (XPM).
\item  Long propagation paths accumulate optical routing loss faster than conventional amplifiers can compensate without introducing amplified spontaneous emission (ASE) corruption.
\item  Thermal diffusion from electronic subsystems destabilizes photonic phase coherence and induces calibration drift.
\end{itemize}

JANUS addresses these constraints through a fully constraint-driven architecture composed of five coupled mechanisms:

\begin{enumerate}
\item One-Hot Optical RNS multiplication for exact low-SNR binary detection.
\item Asymmetric non-volatile PCM 16-Tree Fermat routing networks with the 16th waveguide ($WG_{16}$) extension, requiring only 4 binary stages and 240 switches per multiplier to achieve an optical insertion loss of $1.61~\mathrm{dB}$ with zero static hold power.
\item Single-wavelength coherent spatial fanout removing inter-channel nonlinear interference.
\item High-power 1064 nm passive fan-out paired with Ge/Si $\mathrm{SAC^2M}$ avalanche photodetection, achieving link closure and robust detection margin ($+6.142~\mathrm{dB}$) without intermediate optical amplification or ASE noise accumulation.
\item Spatial CRT reconstruction and JIR-predictive RRNS fault management enabling exact precision scaling with deterministic error detectability.
\end{enumerate}

The resulting system operates entirely within bounded thermal, optical, and electrical regimes:

\begin{equation}
P_{\mathrm{noise}} < P_{\mathrm{signal}} < P_{\mathrm{sat}}, \quad T_{\mathrm{junction}} < T_{\mathrm{critical}}, \quad |\Delta\phi_{\mathrm{drift}}| < \phi_{\mathrm{max}}.
\end{equation}

Rather than pursuing theoretical peak throughput under idealized assumptions, JANUS prioritizes sustained deterministic operation under realistic deployment conditions.

\section{PHYSICAL CONSTRAINT LANDSCAPE}

\subsection{\textbf{Analog Precision Collapse in Photonic Accumulation}}
\label{sec:precision_collapse}

The dominant assumption underlying conventional optical neural networks is that analog optical accumulation can scale proportionally with matrix dimension while maintaining recoverable signal fidelity. This assumption fails under realistic large-scale operating conditions.
For an $N \times N$ photonic matrix multiplication using unsigned 8-bit operands, the worst-case unreduced analog accumulation scales as $S_{\max} = N \times 255 \times 255$, requiring an exact recovery signal-to-noise ratio $\mathrm{SNR}_{\mathrm{req}} = 20\log_{10}(S_{\max})$ \cite{Winzer2006,Saleh2019}:
\begin{equation}
\begin{aligned}
S_{\max,\mathrm{Mini}} &= 2{,}080{,}800 \implies \mathrm{SNR}_{\mathrm{Mini}} \approx 126.4\,\mathrm{dB}, \\
S_{\max,\mathrm{Edge}} &= 4{,}161{,}600 \implies \mathrm{SNR}_{\mathrm{Edge}} \approx 132.4\,\mathrm{dB}, \\
S_{\max,\mathrm{Datacenter}} &= 8{,}323{,}200 \implies \mathrm{SNR}_{\mathrm{Datacenter}} \approx 138.4\,\mathrm{dB}.
\end{aligned}
\label{eq:analog_snr_scaling}
\end{equation}
These requirements exceed realistic photonic operating conditions by several orders of magnitude (as summarized alongside the One-Hot solution in Table~\ref{tab:onehot_transition}).

The contradiction becomes especially severe when detector bandwidth is considered. Because high-resolution requirements---such as the 138 dB analog precision target derived above---cannot be met by standard ADC/DAC architectures operating above 10 GHz , JANUS bypasses these conventional limitations through a specialized spatial encoding approach. As derived in Section~\ref{sec:operating_frequency}, the architecture sustains an operating frequency of 100 GHz; this is critical, as no practical ADC architecture can simultaneously sustain high resolution, high sampling rate, low thermal noise, and acceptable power density within a realistic operating envelope.

Consequently, the analog accumulation paradigm itself becomes the primary scaling failure.

\subsection{\textbf{Failure of Conventional Precision Recovery Techniques}}

Several conventional mitigation strategies attempt to preserve analog photonic scaling.
\begin{itemize}
\item Voltage normalization compresses signal amplitude but does not reduce information density. The required recoverable entropy remains unchanged, meaning the ADC precision burden persists despite lower voltage swing.
\item Temporal bit slicing distributes precision across multiple cycles. If the native throughput is fnative, then the effective throughput becomes

\begin{equation}
f_{\text{eff}} = \frac{f_{\text{native}}}{S}
\end{equation}

where S is the slice factor. For an accumulation of 128 terms requiring 23-bit resolution sliced into 4-bit optical segments, S = 6 cycles per MAC. The effective throughput collapses from 100 GHz to 16.7 GHz, destroying the speed advantage of optics.
\item Spatial interleaving distributes accumulation across multiple physical waveguides. However, the signals must ultimately be coherently combined or individually converted by dedicated ADCs. Slicing across 6 waveguides increases detector and ADC count by 6$\times$, creating an unsustainable area and thermal footprint.
\end{itemize}

Similarly, iterative calibration loops attempt to compensate thermal drift and phase instability through repeated electronic correction.
However, calibration complexity scales with O($N^{2}$) for dense interferometric meshes, making large-scale stable operation thermally impractical. 

Therefore, existing mitigation strategies merely redistribute the precision failure rather than resolving it fundamentally.

\subsection{\textbf{Static Power Explosion in Thermo-Optic Routing Meshes}}
\label{sec:thermo_optic_scaling}

Large-scale silicon photonic accelerators commonly employ electrically reconfigurable optical switching fabrics, including thermo-optic Mach--Zehnder Interferometer (MZI) meshes and other programmable integrated photonic switches, for optical routing and weighting \cite{Miller2013}, \cite{Harris2018} \cite{Zhong2024}.

These architectures require continuous electrical bias to maintain the programmed optical state. Although recent advances have significantly reduced the static tuning power of individual devices, the total static power still scales linearly with the number of programmable switching elements, making power consumption a fundamental challenge for large-scale photonic systems.

Recent state-of-the-art electrically programmable silicon photonic switches have demonstrated static tuning powers as low as

\begin{align}
P_{\mathrm{MZI}}
\approx
0.10~\mathrm{mW}
\label{eq:mzi_static_power}
\end{align}

per active element \cite{Sinatkas2021}.

\paragraph{\textbf{Asymmetric 16-Tree Fermat Core Switch Scaling.}}
In spatial One-Hot modular arithmetic $Y = (A \times B) \pmod{m}$, operand $A$ asserts \textit{exactly one physical waveguide} at any given instant. Because only one input waveguide carries optical power per multiply-accumulate operation, the optical routing engine does not require generalized $N \times N$ arbitrary permutation meshes with deep multi-stage cascaded switches. Instead, the computation is resolved by routing the energized optical pulse through dedicated binary switch trees to the target residue output channel corresponding to weight operand $B$.

JANUS formalizes this principle through the \textbf{Asymmetric 16-Tree Fermat Core} with non-volatile $\mathrm{Sb_2S_3}$ phase-change material (PCM) switching ($0.00~\mathrm{W}$ static hold power). By incorporating the 16th active waveguide ($WG_{16}$) alongside the zero-gated dark reference channel ($WG_0$), the core operates over a 17-channel spatial alphabet ($N_{\mathrm{alphabet}} = 17$):
\begin{equation}
\mathcal{A} = \{ WG_0, WG_1, \dots, WG_{16} \} \longleftrightarrow \{ 0, 1, \dots, 16 \}
\end{equation}
Each of the 16 active waveguides ($WG_1 \dots WG_{16}$) drives an independent 4-stage binary switch tree ($\log_2 16 = 4$ stages). Each binary tree contains $1 + 2 + 4 + 8 = 15$ switches. Thus, each spatial multiplier requires exactly:
\begin{equation}
N_{\mathrm{switch,single}} = 16 \times 15 = 240~\mathrm{switches}
\label{eq:16tree_switches}
\end{equation}
with an ultra-shallow optical path depth of only $S = 4~\mathrm{stages}$.

For a photonic matrix multiplication core operating on $N_{\mathrm{dim}} \times N_{\mathrm{dim}}$ tiles across $N_{\mathrm{tiles}}$ parallel residue channels, the total number of physical switching elements scales as:
\begin{equation}
N_{\mathrm{switch,total}} = N_{\mathrm{tiles}} \times N_{\mathrm{dim}}^2 \times N_{\mathrm{switch,single}}
\label{eq:total_switches}
\end{equation}
Evaluating this across the three deployment configurations demonstrates the unsustainable thermal wall of thermo-optic tuning ($P_{\mathrm{MZI}} \approx 0.10\,\mathrm{mW}$):
\begin{equation}
\begin{aligned}
N_{\mathrm{Mini}} &= 7{,}864{,}320 \implies P_{\mathrm{Mini,TO}} = 786.4\,\mathrm{W}, \\
N_{\mathrm{Edge}} &= 31{,}457{,}280 \implies P_{\mathrm{Edge,TO}} = 3.15\,\mathrm{kW}, \\
N_{\mathrm{Datacenter}} &= 125{,}829{,}120 \implies P_{\mathrm{Datacenter,TO}} = 12.58\,\mathrm{kW}.
\end{aligned}
\label{eq:mzi_heat_wall}
\end{equation}
Because non-volatile $\mathrm{Sb_2S_3}$ PCM switches maintain atomic state with zero holding bias, JANUS slashes this static power to identically $0.00\,\mathrm{W}$ across all tiers.
Non-volatile PCM integration completely eliminates this multi-kilowatt static heat load.

\begin{table*}[htbp]
\centering
\caption{
Thermo-Optic Routing Power Scaling Failure Across Deployment Configurations
}
\label{tab:mzi_failure}
\begin{tabular}{lccc}
\toprule

\textbf{Metric} & \textbf{Mini (32 Tiles)} & \textbf{Edge (32 Tiles)} & \textbf{Datacenter (32 Tiles)} \\
\midrule

Physical Multipliers & $32{,}768$ & $131{,}072$ & $524{,}288$ \\

16-Tree Switch Count ($240$/mult) & $7{,}864{,}320$ & $31{,}457{,}280$ & $125{,}829{,}120$ \\

Thermo-Optic Per-Switch Hold Power & $0.10\text{ mW}$ & $0.10\text{ mW}$ & $0.10\text{ mW}$ \\

Thermo-Optic Static Hold Power & $786.4\text{ W}$ & $3.15\text{ kW}$ & $12.58\text{ kW}$ \\

JANUS PCM Static Hold Power & $\mathbf{0.00\text{ W}}$ & $\mathbf{0.00\text{ W}}$ & $\mathbf{0.00\text{ W}}$ \\

Thermo-Optic Thermal Feasibility & Severely constrained & High runaway risk & Unviable ($>10\text{ kW}$ heat) \\

JANUS PCM Thermal Feasibility & Thermally passive & Thermally passive & Thermally passive \\

Continuous Bias Requirement & Volatile: Yes / PCM: \textbf{No} & Volatile: Yes / PCM: \textbf{No} & Volatile: Yes / PCM: \textbf{No} \\

\bottomrule
\end{tabular}
\end{table*}

The failure mechanism is fundamentally thermodynamic. The generated heat modifies the local refractive index of adjacent waveguides:

\begin{align}
\Delta n = \frac{dn}{dT} \Delta T
\label{eq:thermal_index}
\end{align}

where:

\begin{align}
\frac{dn}{dT} \approx 1.86 \times 10^{-4}~\mathrm{K}^{-1}
\end{align}

for silicon photonics \cite{Saleh2019}, \cite{Incropera2006}.

As routing density increases, lateral heat diffusion produces cross-coupled phase drift across neighboring interferometers. The resulting thermal interference destabilizes the optical transfer matrix and forces continuous recalibration.

Thus, the routing fabric itself becomes the dominant thermal failure mechanism in large-scale photonic systems.

\subsection{\textbf{Nonlinear Failure of Multi-Wavelength Scaling}}
\label{sec:nonlinear_failure}

Many photonic accelerator architectures attempt to improve throughput through Dense Wavelength-Division Multiplexing (DWDM). In these systems, multiple optical carrier frequencies propagate simultaneously through shared waveguides.

While spectrally multiplexed operation increases theoretical bandwidth, it introduces nonlinear optical interference mechanisms that scale superlinearly with optical power density.

The dominant nonlinear effects are:

\begin{itemize}
\item Four-Wave Mixing (FWM)
\item Cross-Phase Modulation (XPM)
\item Self-Phase Modulation (SPM)
\item Spectral Beating Noise
\end{itemize}

In conventional DWDM systems, the nonlinear interference noise power scales approximately as \cite{Boyd2020,Cotter1999}:

\begin{equation}
\sigma_{\mathrm{NL}}^2 \;\propto\; \gamma^2 P^3 L_{\mathrm{eff}}^2 / A_{\mathrm{eff}},
\label{eq:fwm_scaling}
\end{equation}
where $\gamma$ is the nonlinear coefficient, $P$ is the optical power, $L_{\mathrm{eff}}$ is the effective interaction length, and $A_{\mathrm{eff}}$ is the effective mode area. The quadratic dependence on $\gamma$ and $L_{\mathrm{eff}}$ captures the higher-order nature of four-wave mixing, while the $P^3$ term reflects the cubic growth with launch power.

As channel density increases, multiple wavelength combinations satisfy phase-matching conditions, producing ghost frequencies through inter-channel mixing.

The resulting spectral contamination corrupts the recoverability of analog optical information.

Importantly, this failure cannot be completely removed through filtering. Filtering suppresses generated sidebands only after nonlinear mixing has already occurred. Therefore, the underlying epistemic corruption of the optical signal persists.

The problem becomes increasingly severe under high-power amplification, where stronger nonlinear interaction lengths accelerate spectral mixing.

Consequently, wavelength scaling introduces a contradiction:

\textbf{Higher Spectral Density $\Longrightarrow$ Higher Nonlinear Corruption}

This makes large-scale exact arithmetic fundamentally unstable under multi-wavelength analog operation.

\subsection{\textbf{Optical Loss Accumulation and Amplifier Noise}}
\label{sec:loss_accumulation}

Large photonic routing fabrics inevitably accumulate substantial optical propagation loss. The total path loss can be expressed as:

\begin{align}
L_{\mathrm{total}}
&=
L_{\mathrm{mod}}
+
L_{\mathrm{wg}}
+
L_{\mathrm{coupler}}
+
L_{\mathrm{MZI}}
+
L_{\mathrm{routing}}
\label{eq:total_loss}
\end{align}

For practical large-scale systems, component losses typically span $L_{\mathrm{mod}} \approx 3\text{--}6\,\mathrm{dB}$, $L_{\mathrm{wg}} \approx 0.5\text{--}2\,\mathrm{dB}$, $L_{\mathrm{coupler}} \approx 1\text{--}3\,\mathrm{dB}$, and $L_{\mathrm{routing}} \approx 15\text{--}20\,\mathrm{dB}$, yielding cumulative uncompensated attenuation $L_{\mathrm{total}} \approx 30\text{--}50\,\mathrm{dB}$.


Without compensation, optical power decays exponentially along the propagation path:

\begin{align}
P(z) = P_0 e^{-\alpha z}
\label{eq:loss_exponential}
\end{align}

where: $\alpha$ is the effective attenuation coefficient.

Conventional semiconductor optical amplifiers (SOAs) attempt to restore signal amplitude through stimulated emission. However, SOAs introduce Amplified Spontaneous Emission (ASE) noise \cite{Winzer2006,Saleh2019}:

\begin{align}
\sigma^2_{\mathrm{ASE}} \propto Gn_{sp}B
\label{eq:ase_noise}
\end{align}

where: $G$ is amplifier gain, $n_{sp}$ is the spontaneous emission factor, and $B$ is bandwidth.

The contradiction is therefore:

\textbf{Higher Gain $\Longrightarrow$ Higher ASE Corruption}

Thus, conventional amplification preserves optical amplitude while simultaneously degrading informational integrity.

For exact arithmetic systems, this failure is especially severe because even small analog corruption propagates catastrophically through high-dimensional reconstruction.

\subsection{\textbf{Thermal Drift and Calibration Instability}}
\label{sec:thermal_drift}

Photonic interferometric systems are intrinsically sensitive to thermal variation because optical phase depends directly upon refractive index and propagation length.

The accumulated optical phase is:

\begin{align}
\phi = \frac{2\pi nL}{\lambda}
\label{eq:phase_equation}
\end{align}

Thermal perturbation therefore produces phase drift:

\begin{align}
\Delta \phi = \frac{2\pi L}{\lambda} \frac{dn}{dT} \Delta T
\label{eq:phase_drift}
\end{align}

As routing density increases, neighboring thermal sources generate dynamic phase instability throughout the photonic mesh.

The resulting transfer matrix becomes time-varying:

\begin{align}
H(t) \neq H(t+\Delta t)
\label{eq:time_variation}
\end{align}

Conventional architectures compensate this instability through active calibration loops involving:

\begin{itemize}
\item thermal tuning,
\item phase locking,
\item electronic correction,
\item iterative optimization,
\item feedback control.
\end{itemize}

However, calibration overhead scales poorly with mesh density \cite{Miller2013}, \cite{Hamerly2022}:

$C_{\mathrm{calibration}} \propto N^2$

At large dimensions, calibration itself becomes a dominant energy and latency cost.

Furthermore, continuously recalibrated systems do not operate within a stable thermodynamic equilibrium. Instead, they remain permanently trapped in transient correction states. Thus, large-scale photonic accelerators face a second-order contradiction:



This contradiction prevents sustained deterministic operation under realistic deployment conditions.

\subsection{\textbf{Constraint-Driven Architectural Reformulation}}

The preceding constraints collectively imply that large-scale photonic exact arithmetic cannot emerge from incremental improvement of conventional analog architectures.

The failure mechanisms are structural rather than implementation-level.

Specifically:

\begin{itemize}
\item  Analog accumulation fails due to Shannon-limited precision collapse.
\item  Thermo-optic routing fails due to unsustainable static power density.
\item  Multi-wavelength scaling fails due to nonlinear spectral corruption.
\item  Conventional amplification fails due to ASE-induced epistemic loss.
\item  Dense interferometric meshes fail due to thermal calibration instability.
\end{itemize}

Consequently, JANUS adopts a fundamentally different computational philosophy. Rather than attempting to improve analog reconstruction fidelity, JANUS reformulates optical computation into a bounded spatial routing problem.

The resulting architecture enforces $P_{\mathrm{noise}} < P_{\mathrm{signal}} < P_{\mathrm{sat}}$, $T_{\mathrm{junction}} < T_{\mathrm{critical}}$, and $|\Delta\phi_{\mathrm{drift}}| < \phi_{\mathrm{max}}$ across all operating regimes.

This bounded-operation methodology intentionally sacrifices theoretical peak analog density in favor of sustained exact arithmetic under realistic physical constraints.

\section{CONSTRAINT-DRIVEN ARCHITECTURAL RESPONSE}


\subsection{\textbf{Architectural Overview}}

JANUS resolves the scaling contradictions identified in the previous section through a tightly coupled set of architectural mechanisms, each directly derived from a specific physical failure mode.

The resulting system is not an incremental extension of conventional analog photonic computing. Instead, it is a reformulation of photonic computation around bounded spatial exact arithmetic.

The architecture consists of six tightly integrated subsystems:

\begin{enumerate}
\item One-Hot Optical Residue Number System multiplication,
\item Asymmetric PCM 16-Tree Fermat routing fabrics with $WG_{16}$ extension,
\item High-power 1064 nm passive fan-out and APD detection,
\item Single-wavelength coherent spatial fanout,
\item Spatial CRT precision reconstruction,
\item Thermally decoupled 3D Z-axis integration.
\end{enumerate}

Each subsystem exists because a conventional alternative becomes physically unstable at large scale.
The resulting architecture operates under a constraint-first design principle:
\textbf{Physical Stability > Theoretical Peak Density}

Unlike conventional photonic accelerators that rely upon continuous calibration and optimistic analog assumptions, JANUS constrains all thermal, optical, and electrical states within bounded operating regions.

This bounded-operation philosophy forms the foundational principle of the entire architecture.

\paragraph{\textbf{Demonstrated Components and Forward-Looking Integration.}}

Several device-level components discussed in this work have been independently demonstrated in literature, including 105~GHz Ge/Si $\mathrm{SAC^2M}$ avalanche photodetectors \cite{Chen2025}, non-volatile $\mathrm{Sb_2S_3}$ PCM switching in silicon photonic waveguides utilizing graphene micro-heaters \cite{Delaney2020,Rios2021}, and silicon photonic routing meshes. However, the full JANUS architecture as described---including the high-power passive fan-out and APD link closure, multi-million-switch non-volatile 16-Tree Fermat fabric ($125.8\text{ M}$ switches in Datacenter, $7.86\text{ M}$ in 32-tile Mini, and $3.93\text{ M}$ in the 16-tile Model~1A prototype), complete stacked integration, and multi-tile datacenter configuration---remains a forward-looking systems integration proposal. This paper presents the architectural and physical motivation for such integration rather than a fabrication-validated implementation.

\begin{figure*}[!t]
\centering
\begin{tikzpicture}[
    % ---- global font ----
    every node/.style={font=\small\sffamily},
    % ---- node styles ----
    input/.style={
        draw, rounded corners=3pt,
        fill=gray!12, minimum width=2.6cm, minimum height=0.72cm,
        align=center, thick
    },
    stage/.style={
        draw, rounded corners=3pt,
        fill=blue!8, minimum width=2.8cm, minimum height=0.80cm,
        align=center, thick
    },
    optical/.style={
        draw, rounded corners=3pt,
        fill=cyan!10, minimum width=2.8cm, minimum height=0.80cm,
        align=center, thick
    },
    cmos/.style={
        draw, rounded corners=3pt,
        fill=orange!10, minimum width=2.8cm, minimum height=0.80cm,
        align=center, thick
    },
    output/.style={
        draw, rounded corners=3pt,
        fill=green!10, minimum width=2.6cm, minimum height=0.72cm,
        align=center, thick
    },
    fault/.style={
        draw=red!70, rounded corners=3pt,
        fill=red!6, minimum width=2.4cm, minimum height=0.68cm,
        align=center, thick, dashed
    },
    annot/.style={
        font=\scriptsize\sffamily\itshape, text=gray!70, align=center
    },
    % ---- arrow styles ----
    arr/.style={-{Stealth[scale=1.1]}, thick},
    arropt/.style={-{Stealth[scale=1.1]}, thick, color=cyan!60!black},
    arrcmos/.style={-{Stealth[scale=1.1]}, thick, color=orange!70!black},
    arrfault/.style={-{Stealth[scale=1.0]}, thick, color=red!60, dashed},
    % ---- lane separation ----
    node distance=0.55cm and 0.0cm
]

% =============================================================
%  ROW 0  -- Input
% =============================================================
\node[input] (IN) {Integer Input $X$\\(64-bit operand)};

% =============================================================
%  ROW 1  -- RNS Decomposition  (CMOS pre-processing)
% =============================================================
\node[cmos, below=0.65cm of IN] (RNS)
    {RNS Decomposition\\$x_i = X \bmod m_i$};

% Arrow: Integer Input X -> RNS Decomposition
\draw[arr] (IN.south) -- (RNS.north);

% fan-out annotation
\node[annot, below=0.20cm of RNS] (fanannot)
    {$k$ independent residue channels\\$x_i \in [0,\,255]$ each};

% =============================================================
%  ROW 2  -- Residue tiles  (three shown, dots for rest)
%  Laid out horizontally below the RNS block
% =============================================================
\node[optical, below=1.05cm of fanannot, xshift=-3.6cm] (T1)
    {Residue Tile $1$\\mod $m_1$};
\node[optical, below=1.05cm of fanannot]               (T2)
    {Residue Tile $2$\\mod $m_2$};
\node[optical, below=1.05cm of fanannot, xshift=+3.6cm] (Tk)
    {Residue Tile $k$\\mod $m_k$};

% dots between tiles
\node[annot, below=1.05cm of fanannot, xshift=+1.85cm] (dots1) {$\cdots$};
\node[annot, below=1.05cm of fanannot, xshift=-1.85cm] (dots2) {$\cdots$};

% arrows from RNS to each tile
\draw[arrcmos] (RNS.south) -- ++(0,-0.30) -| (T1.north);
\draw[arrcmos] (RNS.south) -- ++(0,-0.30) -| (T2.north);
\draw[arrcmos] (RNS.south) -- ++(0,-0.30) -| (Tk.north);

% =============================================================
%  ROW 3  -- Inside each tile: One-Hot encoding + fanout
% =============================================================
\node[optical, below=0.60cm of T1] (OH1)
    {One-Hot Encoding\\$x_i \to e_{x_i}$\\17 waveguides};
\node[optical, below=0.60cm of T2] (OH2)
    {One-Hot Encoding\\$x_i \to e_{x_i}$\\17 waveguides};
\node[optical, below=0.60cm of Tk] (OHk)
    {One-Hot Encoding\\$x_i \to e_{x_i}$\\17 waveguides};

\draw[arropt] (T1.south) -- (OH1.north);
\draw[arropt] (T2.south) -- (OH2.north);
\draw[arropt] (Tk.south) -- (OHk.north);

% single-wavelength fanout note inside tile 2
\node[annot, right=4 cm of OH2] (swl)
    {single-$\lambda$\\coherent\\fanout};

% =============================================================
%  ROW 4  -- PCM 16-Tree Fermat routing fabric
% =============================================================
\node[optical, below=0.60cm of OH1] (B1)
    {PCM 16-Tree Fermat\\Routing Fabric\\($\log_2 16 = 4$ stages)};
\node[optical, below=0.60cm of OH2] (B2)
    {PCM 16-Tree Fermat\\Routing Fabric\\($\log_2 16 = 4$ stages)};
\node[optical, below=0.60cm of OHk] (Bk)
    {PCM 16-Tree Fermat\\Routing Fabric\\($\log_2 16 = 4$ stages)};

\draw[arropt] (OH1.south) -- (B1.north);
\draw[arropt] (OH2.south) -- (B2.north);
\draw[arropt] (OHk.south) -- (Bk.north);

% Passive link closure annotation
\node[annot, right=4cm of B2] (psa)
    {Passive 1064 nm\\Fan-Out + Ge/Si\\$\mathrm{SAC^2M}$ APDs};

% =============================================================
%  ROW 5  -- Binary photodetection
% =============================================================
\node[optical, below=0.60cm of B1] (D1)
    {Ge/Si $\mathrm{SAC^2M}$ APD\\$D_k{=}1$ if $P_\text{sig}{>}P_\text{sens}$\\residue $z_1$};
\node[optical, below=0.60cm of B2] (D2)
    {Ge/Si $\mathrm{SAC^2M}$ APD\\$D_k{=}1$ if $P_\text{sig}{>}P_\text{sens}$\\residue $z_2$};
\node[optical, below=0.60cm of Bk] (Dk)
    {Ge/Si $\mathrm{SAC^2M}$ APD\\$D_k{=}1$ if $P_\text{sig}{>}P_\text{sens}$\\residue $z_k$};

\draw[arropt] (B1.south) -- (D1.north);
\draw[arropt] (B2.south) -- (D2.north);
\draw[arropt] (Bk.south) -- (Dk.north);

% =============================================================
%  ROW 6  -- CRT Reconstruction  (CMOS back-end)
% =============================================================
\node[cmos,
      below=1.10cm of D2,
      minimum width=9.8cm] (CRT)
    {CRT Reconstruction (CMOS back-end):\quad
     $X = \displaystyle\sum_{i=1}^{k} z_i\, M_i\, N_i \bmod M$};

\draw[arrcmos] (D1.south) -- ++(0,-0.30) -| (CRT.north);
\draw[arrcmos] (D2.south) -- (CRT.north);
\draw[arrcmos] (Dk.south) -- ++(0,-0.30) -| (CRT.north);

% =============================================================
%  ROW 7  -- CRT consistency check
% =============================================================
\node[cmos, below=0.60cm of CRT] (LH)
    {CRT Consistency Check\\JIR residue cross-verification};

\draw[arrcmos] (CRT.south) -- (LH.north);

% =============================================================
%  ROW 8  -- Valid output  OR  Fault signal
% =============================================================
\node[output,  below=0.65cm of LH, xshift=-2.0cm] (OUT)
    {Valid Output $X$\\(64-bit exact result)};
\node[fault,   below=0.65cm of LH, xshift=+2.4cm] (FAULT)
    {Explicit Fault\\(CRT invalidation)};

\draw[arrcmos]  (LH.south) -- ++(0,-0.22) -| (OUT.north);
\draw[arrfault] (LH.south) -- ++(0,-0.22) -| (FAULT.north);

% labels on the two output branches
\node[annot, below=0.08cm of LH, xshift=-1.05cm] {\scriptsize consistent};
\node[annot, below=0.08cm of LH, xshift=+1.30cm] {\scriptsize inconsistent};

% =============================================================
%  LEFT-SIDE LAYER LABELS  (Z-axis stack annotation)
% =============================================================
\node[annot, left=0.5cm of RNS,  xshift=-0.3cm]
    {\textbf{CMOS}\\layer};

\node[annot, left=0.5cm of OH1,  xshift=-0.3cm]
    {\textbf{SiPh}\\layer};

\node[annot, left=0.5cm of B1,   xshift=-0.3cm]
    {\textbf{SiPh}\\(routing)};

\node[annot, left=0.5cm of D1,   xshift=-0.3cm]
    {\textbf{Detector}\\layer};

\node[annot, left=0.5cm of CRT,  xshift=-0.3cm]
    {\textbf{CMOS}\\back-end};

% =============================================================
%  RIGHT-SIDE  non-volatile PCM note
% =============================================================
\node[annot, right=2 cm of B2, xshift=+1.3cm, yshift=1 cm]
    {non-volatile state:\\zero static hold power};

% =============================================================
%  SURROUNDING BOUNDING BOX  -- Photonic domain
% =============================================================
\begin{pgfonlayer}{background}
  % shade the optical region
\end{pgfonlayer}

\end{tikzpicture}

\caption{JANUS end-to-end data and computation flow. An input word is decomposed into $k$ independent residues by the CMOS front-end. Each residue is encoded spatially as a one-hot vector across optical waveguides within its dedicated optical tile. A coherent single-wavelength fanout distributes light; a non-volatile Asymmetric 16-Tree Fermat PCM fabric performs deterministic spatial permutation corresponding to modular multiplication across 4 stages ($\log_2 16 = 4$, $240$ switches/mult), with the 16th waveguide ($WG_{16}$) natively unlocking $\mathbb{Z}_{17}$ state efficiency ($100\%$) and Radix-16 $\mathbb{Z}_{257}$ division-free modular reduction. High-power 1064 nm passive fan-out and Ge/Si $\mathrm{SAC^2M}$ avalanche photodetectors provide link closure without intermediate optical amplification. Binary photodetectors read the active waveguide position, producing the modular product $z_i$. The CMOS back-end reconstructs the exact multi-precision result via the Chinese Remainder Theorem and applies JIR/RRNS consistency checking.}
\label{fig:janus_dataflow}
\end{figure*}

% ============================================================
%  END OF DIAGRAM BLOCK
% ============================================================

\subsection{\textbf{One-Hot Optical Residue Number System Multiplication}}

The central architectural transition within JANUS is the elimination of analog optical accumulation. Instead of encoding numerical information into continuously varying optical amplitudes, JANUS reformulates arithmetic as a spatial routing problem using One-Hot Optical Residue Number System (RNS) encoding \cite{Omosanya2021}\cite{Huang1979}\cite{Sorger2018}.

This transition is motivated by the observation that analog amplitude reconstruction constitutes the dominant precision bottleneck in large-scale photonic computation.

As established in Section~\ref{sec:precision_collapse}, conventional analog optical accumulation requires an unachievable $>138\,\mathrm{dB}$ SNR at multi-GHz frequencies. JANUS eliminates this analog precision bottleneck by reformulating arithmetic into a deterministic spatial localization problem over One-Hot optical residue channels.

\subsubsection{\textbf{Transition from Amplitude Reconstruction to Spatial Localization}}
 
JANUS resolves this contradiction by abandoning analog amplitude representation entirely.Instead of encoding values through optical intensity, JANUS encodes values spatially. Each residue value activates exactly one optical waveguide among the spatial channels (e.g., 17 channels in the Asymmetric 16-Tree Fermat Core):

$x \mapsto e_x$

where:
\begin{equation}
e_x = (0, 0, \dots, 1, \dots, 0)
\end{equation}

is a one-hot spatial basis vector.

Thus, arithmetic state is represented through waveguide position rather than analog optical amplitude.

Optical multiplication becomes deterministic spatial routing:

\begin{equation}
(a, b) \rightarrow (a \cdot b) \mod p
\label{eq:rns_multiplication}
\end{equation}

where: p is the local residue modulus.

The photodetector therefore no longer reconstructs analog magnitude. Instead, it performs binary localization:

\begin{align}
D_k
=
\begin{cases}
1,
&
\text{light present in channel } k
\\
0,
&
\text{otherwise}
\end{cases}
\label{eq:binary_detection}
\end{align}

This transforms the computational problem from analog amplitude estimation to binary spatial localization.

The required detector condition therefore collapses to:

\begin{center}
\text{$P_{signal} > P_{noise}$}
\label{eq:binary_condition}
\end{center}

which corresponds approximately to:

\begin{center}
\text{$\mathrm{SNR}_{\mathrm{binary}} > 0~\mathrm{dB}$}
\end{center}

\begin{equation}
\left(
P_{\mathrm{signal}} > P_{\mathrm{sens,practical}}
\approx -23.2~\mathrm{dBm}
\right)
\end{equation}

The Ge/Si $\mathrm{SAC^2M}$ avalanche photodetector operating at gain $M = 7$ possesses a gain--bandwidth product of 441 GHz with an inherent 3 dB receiver bandwidth of 105 GHz, comfortably supporting the 100 GHz system clock rate. This bandwidth filtering combined with a receiverless charge-integration front end and an ILO comb-referenced clock (50 fs rms jitter) yields a practical binary receiver sensitivity of $P_{\mathrm{sens,practical}} \approx -23.20\text{ dBm}$ (for a reconciled target bit error rate of $\le 10^{-36}$, as derived in Section~\ref{sec:passive_fanout_apd}). Consequently, the delivered optical signal intensity across the passive 1 : $2^{18}$ distribution tree reaches $-12.70\text{ dBm}$ through the 4-stage 16-Tree Fermat Core, providing an abundant binary detection margin of $+10.50\text{ dB}$ ($+6.142\text{ dB}$ under strict foundry sign-off constraints).

Thus, the JANUS architecture transforms an analog precision requirement of approximately 138 dB into a robust binary detection problem with substantial receiver margin.

\begin{equation}
P_{\mathrm{th}} \approx -23.20~\mathrm{dBm}
\end{equation}

By utilizing highly sensitive avalanche photodetectors (APDs), JANUS guarantees that even after maximal passive splitting, the terminal optical power exceeds the detection threshold without requiring intermediate re-amplification (see Section~\ref{sec:passive_fanout_apd}).

\begin{table}[htbp]
\centering
\footnotesize
\setlength{\tabcolsep}{2.5pt}
\renewcommand{\arraystretch}{1.15}
\caption{TRANSITION FROM ANALOG ACCUMULATION TO ONE-HOT SPATIAL DETECTION}
\label{tab:onehot_transition}
\begin{tabularx}{\columnwidth}{@{}Xccc@{}}
\toprule
\textbf{Metric} & \textbf{Mini} & \textbf{Edge} & \textbf{Datacenter} \\
\midrule
Max. Accumulation & $2{,}080{,}800$ & $4{,}161{,}600$ & $8{,}323{,}200$ \\
Required Analog SNR & $126.4\,\mathrm{dB}$ & $132.4\,\mathrm{dB}$ & $138.4\,\mathrm{dB}$ \\
Sensitivity ($P_{\mathrm{th}}$) & $-23.20\,\mathrm{dBm}$ & $-23.20\,\mathrm{dBm}$ & $-23.20\,\mathrm{dBm}$ \\
ADC Requirement & None (0-bit) & None (0-bit) & None (0-bit) \\
Detection Scheme & Binary OOK & Binary OOK & Binary OOK \\
Operational Stability & Bounded & Bounded & Bounded \\
\bottomrule
\end{tabularx}
\end{table}

\subsubsection{\textbf{One-Hot Residue Arithmetic}}

 JANUS performs arithmetic independently within residue domains.
For a modulus set:
\begin{equation}
\mathcal{M} = \{m_1, m_2, \dots, m_k\}
\label{eq:modulus_set}
\end{equation}

a large integer $X$ is decomposed into residue components:
\begin{equation}
X \rightarrow (x_1, x_2, \dots, x_k)
\end{equation}

where:
\begin{equation}
x_i = X \mod m_i
\label{eq:residue_projection}
\end{equation}

Each residue channel operates independently:
\begin{equation}
z_i = (x_iy_i) \mod m_i
\label{eq:modular_product}
\end{equation}

Importantly, no analog carry propagation exists between channels. This eliminates the principal scaling bottleneck of conventional binary arithmetic \cite{Omosanya2021}, \cite{Mohan2016}:

$\text{Carry Propagation Delay} \to 0$

because each residue tile computes independently.
The resulting architecture therefore scales spatially rather than temporally.

\subsubsection{\textbf{Optical Routing Interpretation of Multiplication}}
Within JANUS, multiplication is implemented physically as deterministic waveguide permutation.

For modulus m the routing function becomes:
\begin{equation}
f (a, b) = (a \cdot b) \mod m
\end{equation}

The optical mesh therefore behaves as a programmable permutation tensor.

Unlike analog MAC systems, no weighted accumulation occurs. Optical energy is conserved spatially:

\begin{equation}
\sum_{i=1}^k P_i = P_{\mathrm{total}}
\end{equation}

Modular multiplication within channel $m_i$ takes the form:

\begin{equation}
z_i = (x_i \cdot w_i) \bmod m_i
\end{equation}

Because $w_i$ is known during matrix execution, multiplication by a constant weight $w_i$ defines a fixed bijective permutation on the residue elements:

\begin{equation}
\pi_{w_i}: x_i \mapsto (x_i \cdot w_i) \bmod m_i
\end{equation}

Optically, a permutation of $m_i$ channels is realized without arithmetic computation: it corresponds strictly to a spatial permutation of physical waveguides. The optical signal propagating along the input waveguide $x_i$ is routed deterministically to the output waveguide $z_i$.

\subsubsection{\textbf{Detection Margin and False Positive Suppression}}

Binary detection operates under the hypothesis test:

\begin{center}
\text{$H_0$: No photon present (dark channel)}
\end{center}

\begin{center}
\text{$H_1$: Photon present (active channel)}
\end{center}

The decision threshold $P_{\mathrm{th}}$ is set at the receiver sensitivity level:

\begin{equation}
P_{\mathrm{th}} = P_{\mathrm{sens,practical}} \approx -23.20\,\text{dBm} \approx 4.79\,\mu\text{W}
\end{equation}

A detection event is registered if and only if the received optical power $P_{\mathrm{rec}}$ satisfies:

\begin{equation}
P_{\mathrm{rec}} \ge P_{\mathrm{th}}
\end{equation}

The detection margin is defined as:

\begin{equation}
M = P_{\text{signal}} - P_{\text{noise}}
\end{equation}

For the nominal passive fan-out routing path, the optical chain proceeds as follows: the master 1064 nm Yb-fiber laser ($P_{\text{laser}} = +50.0\text{ dBm}$) is passively split across the 1 : $2^{18}$ optical distribution tree (ideal splitting loss 54.19 dB, plus excess path loss comprising 5.40 dB MMI excess loss, $1.61\,\text{dB}$ 4-stage 16-Tree routing loss, and 1.50 dB propagation and coupling penalties, yielding $L_{\text{total}} = 62.70\text{ dB}$). Thus, the delivered optical power at each terminal Ge/Si APD receiver reaches:

\begin{equation}
P_{\text{signal}} \equiv P_{\text{det}} \approx -12.70\text{ dBm} \approx 53.7\,\mu\text{W}.
\end{equation}

As derived in Section~\ref{sec:passive_fanout_apd}, the Ge/Si $\mathrm{SAC^2M}$ avalanche photodetector operating at gain $M = 7$ with a receiverless charge-integration front end and an ILO comb-referenced clock (50 fs rms jitter) achieves a practical binary OOK detection sensitivity for a target bit error rate of $\le 10^{-36}$ (reconciled integrate-and-dump $Q \approx 12.72$) of:

\begin{equation}
P_{\text{sens,practical}} \approx -23.20\text{ dBm}.
\end{equation}

Therefore, the binary detection margin across the full passive fan-out and routing network is:

\begin{equation}
\text{Margin} = (-12.70\text{ dBm}) - (-23.20\text{ dBm}) = +10.50\text{ dB} \quad (+6.142\text{ dB under strict sign-off}).
\end{equation}

Because one-hot encoding guarantees that exactly one waveguide per multiplier contains optical power (with 16 unselected dark waveguides in the 17-channel Asymmetric 16-Tree Fermat Core), the false positive probability across all dark channels is bounded by the bit error rate:

\begin{equation}
P_{\mathrm{FP}} \le 16 \times \mathrm{BER} = 16 \times 10^{-18} = 1.60 \times 10^{-17}
\end{equation}
(suppressed to $\le 1.60 \times 10^{-35}$ under the reconciled integrate-and-dump $\mathrm{BER} \le 10^{-36}$ regime).

This confirms that binary spatial detection operates with extreme fidelity across the entire passive fan-out without requiring intermediate optical amplification.

\subsubsection{\textbf{Architectural Consequences}}

The spatial one-hot representation establishes several fundamental architectural guarantees:

\begin{enumerate}
\item Precision decoupling: Arithmetic correctness depends solely on spatial routing path validity, not on analog amplitude fidelity.
\item Noise immunity: Signal attenuation reduces link margin but introduces no numerical error as long as delivered power exceeds $P_{\mathrm{th}}$.
\item Dynamic range scalability: High-precision computation is achieved through parallel residue channels rather than higher analog SNR.
\end{enumerate}

\subsection{\textbf{Asymmetric Non-Volatile PCM 16-Tree Fermat Routing Fabric}}

To implement spatial permutations $\pi_{w_i}$ efficiently without static power dissipation, JANUS incorporates non-volatile phase-change material (PCM) switches arranged in an Asymmetric 16-Tree Fermat Core topology \cite{Fang2022}\cite{Zhang2021}\cite{Wuttig2017}\cite{Sinatkas2021}. Rather than deploying an arbitrary non-blocking permutation mesh, the 16-Tree structure exploits One-Hot spatial sparsity: each active input waveguide routes light independently through a dedicated 4-stage binary tree with zero static power dissipation.

As quantified in Section~\ref{sec:thermo_optic_scaling}, continuously biased thermo-optic MZI switches dissipate multi-kilowatts of static heater power at scale ($12.58\,\mathrm{kW}$ in Datacenter). JANUS eliminates static hold power entirely by replacing thermo-optic heaters with non-volatile $\mathrm{Sb_2S_3}$ phase-change material (PCM) switches ($P_{\mathrm{hold,PCM}} = 0\,\mathrm{W}$).

\subsubsection{\textbf{PCM-Based Non-Volatile Optical Routing}}

JANUS eliminates continuous static hold power by replacing thermo-optic MZIs with non-volatile phase-change material switches \cite{Xu2019,Zhang2019,Delaney2020,Rios2021}. Phase-change materials exhibit large refractive index contrast between amorphous and crystalline structural phases. However, legacy tellurides (e.g., GST-225, GST-467) and intermediate alloys (e.g., GSST) suffer from prohibitive resonant interband absorption at the 1064 nm operating wavelength. 

The choice of optical wavelength ($\lambda_0 = 1064\,\mathrm{nm}$, $h\nu \approx 1.165\,\mathrm{eV}$) and the PCM bandgap ($E_g$) must avoid massive resonant interband absorption. For materials like GST-467 or GSST with $E_g \approx 0.50\text{--}0.75\,\mathrm{eV}$, incoming photons exceed the bandgap ($1.165\,\mathrm{eV} > E_g$), resulting in an extremely large extinction coefficient ($\kappa \gg 0.1$). This causes multi-stage fabric loss to exceed 36.75 dB for GST-467 and 17.25 dB for GSST, completely blocking or severely degrading the optical link.

JANUS selects the wide-bandgap chalcogenide $\mathrm{Sb_2S_3}$ (Antimony Trisulfide) \cite{Delaney2020}. Because its electronic bandgap $E_g \approx 1.72\,\mathrm{eV}$ is strictly greater than the photon energy ($1.72\,\mathrm{eV} > 1.165\,\mathrm{eV}$), single-photon interband transitions are forbidden. The material operates in the sub-bandgap transparency window, yielding an extinction coefficient $\kappa < 8.0 \times 10^{-5}$ while retaining large refractive index contrast $\Delta n \approx 0.60$:

\begin{equation}
n_{\mathrm{cryst}} \approx 3.30, \quad n_{\mathrm{amorph}} \approx 2.70 \implies \Delta n = 0.60
\end{equation}

After programming, the optical state persists without continuous power. Thus:

\begin{equation}
P_{\mathrm{hold,PCM}} = 0~\mathrm{W}
\end{equation}

Energy is consumed strictly during weight reconfiguration:

\begin{equation}
E_{\mathrm{reconfig}} = N_{\mathrm{reconfig}} \times E_{\mathrm{switch}}
\end{equation}

where $E_{\mathrm{switch}} \approx 4.2\text{ pJ}$ per cell.

% [Table 4 integrated into Table 1]


\subsubsection{\textbf{$\mathrm{Sb_2S_3}$ Phase-Change Material and Graphene Micro-Heaters}}

To achieve long-term cycling endurance ($> 10^8$ cycles) and minimize optical switching energy across millions of non-volatile 16-Tree Fermat routing switches, JANUS adopts $\mathrm{Sb_2S_3}$ co-integrated with monolayer graphene micro-heaters \cite{Rios2021}.

\paragraph{\textbf{Sub-Bandgap Transparency at 1064 nm \& Sub-0.5 dB Physical Roadmap}}
Legacy plasmonic phase-change switches predominantly utilized $\mathrm{Ge_2Sb_2Te_5}$ (GST-225) or $\mathrm{Ge_4Sb_6Te_7}$ (GST-467). However, these narrow-bandgap materials induce prohibitive interband absorption loss at near-infrared wavelengths. While mainstream integrated photonics research has historically concentrated on telecom wavelengths ($1550~\mathrm{nm}$) due to telecommunications transceiver heritage \cite{Zhou2024,Bao2025}, JANUS deliberately targets $\lambda_0 = 1064\,\mathrm{nm}$ to leverage the massive wall-plug efficiency of Yb-fiber lasers ($75\%$ WPE vs $15-20\%$ for telecom InP DFB lasers).

At $\lambda_0 = 1064\,\mathrm{nm}$, the photon energy ($h\nu = 1.165\,\mathrm{eV}$) is deep within the sub-bandgap transparency window of $\mathrm{Sb_2S_3}$ ($E_g \approx 1.72\,\mathrm{eV}$, $\lambda_g \approx 720\,\mathrm{nm}$), ensuring near-zero intrinsic material absorption ($\kappa \approx 10^{-5}$ in both amorphous and crystalline phases \cite{Delaney2020,Zhang2023MDM}). Non-volatile pulsed switching at $1064-1080\,\mathrm{nm}$ has experimentally demonstrated high optical contrast ($22\,\mathrm{dB}$) under $10\,\mathrm{ps}$ pulses \cite{PerezFrances2023}. 

Because intrinsic material absorption is negligible ($\kappa < 10^{-4}$), total insertion loss is determined by extrinsic fabrication scattering and modal engineering. JANUS achieves sub-$0.50\,\mathrm{dB}$ (and down to $0.10\,\mathrm{dB}$) insertion loss through four physical fabrication and design techniques:
\begin{enumerate}
\item \textbf{Stoichiometric PVD Sputtering:} Ultra-pure stoichiometric magnetron sputtering avoids sulfur deficiency and suppresses phase-transition micro-voids.
\item \textbf{Conformal ALD Passivation:} Atomic layer deposition of conformal $\mathrm{Al_2O_3}$ or $\mathrm{Si_3N_4}$ encapsulation smooths etched waveguide sidewall roughness and suppresses scattering.
\item \textbf{Analytical Phase De-Coupling Geometry:} In the asymmetric $2 \times 2$ directional coupler, complete cancellation of cross-port leakage occurs at the first de-coupling cancellation node $S \cdot L = \pi$ (where $S = \sqrt{\kappa_0^2 + (\Delta\beta/2)^2}$), satisfying:
\begin{equation}
\Gamma \cdot L = \frac{\sqrt{3}\lambda_0}{2\Delta n} \approx \frac{0.866 \times 1064\,\mathrm{nm}}{0.60} \approx 1.536\,\mu\mathrm{m}
\end{equation}
Designing the switch at $L = 60.0\,\mu\mathrm{m}$ with modal overlap $\Gamma = 2.56\%$ satisfies this condition, yielding nominal simulated insertion losses of $S_{31} = -0.010\,\mathrm{dB}$ (amorphous cross-state) and $S_{21} = -0.096\,\mathrm{dB}$ (crystalline bar-state) with crosstalk suppressed below $-74.5\,\mathrm{dB}$.
\item \textbf{Graphene Micro-Heater Actuation:} Monolayer graphene heaters provide localized $4.2\,\mathrm{pJ/cell}$ electro-thermal switching without introducing metallic optical absorption into the waveguide mode.
\end{enumerate}

\paragraph{\textbf{Optical Link Margin Enhancement}}
While the conservative baseline sign-off ceiling adopts $0.50\,\mathrm{dB/cell}$, the 4-stage Asymmetric 16-Tree Fermat Core achieves a total fabric insertion loss of just $1.61\,\mathrm{dB}$ (with measured amorphous cell loss of $0.2627\,\mathrm{dB}$). This provides $+5.89\,\mathrm{dB}$ of optical power headroom above design requirements, elevating the verified operating link margin to $\mathbf{+6.142\,\mathrm{dB}}$ and ensuring link closure across high-fanout distribution trees.

\paragraph{\textbf{Graphene Micro-Heater Co-Integration}}
To switch $\mathrm{Sb_2S_3}$ between its amorphous (CROSS, $n=2.70$) and crystalline (BAR, $n=3.30$) states without free-space optical writing lasers, JANUS integrates monolayer graphene micro-heaters directly above the $\mathrm{Sb_2S_3}$ patches.

The ultra-low heat capacity ($< 0.1\text{ pJ/K}$) of the 1.0 nm graphene heater enables ultra-fast thermal switching with a $4.2\text{ pJ/cell}$ energy footprint. 

\begin{itemize}
\item \textbf{Crystallization (SET Pulse):} A low-voltage electrical pulse heats the $\mathrm{Sb_2S_3}$ patch to $T_c \approx 200\text{--}220^\circ\mathrm{C}$ for $\approx 100\,\mathrm{ns}$.
\item \textbf{Amorphization (RESET Pulse):} A short, high-intensity pulse heats the patch above its melting point ($T_m \approx 500\text{--}550^\circ\mathrm{C}$) within $\approx 15\,\mathrm{ns}$, followed by ultra-fast thermal quenching ($> 10^{10}\,\mathrm{K/s}$) through the Si substrate.
\end{itemize}

Once switched, the non-volatile atomic structure remains locked indefinitely ($> 10\text{ years}$ data retention), requiring 0 W holding power.

\subsubsection{\textbf{Asymmetric 16-Tree Fermat Core Architecture}}
\label{sec:16tree_arch}

\paragraph{\textbf{One-Hot Spatial Modular Routing Principle:}}
In One-Hot spatial Residue Number System (RNS) multiplication $Y = (A \times B) \pmod{m}$, input operand $A$ asserts \textit{exactly one spatial waveguide} during any given clock cycle. Unlike telecommunication optical routing fabrics that must simultaneously permute $N$ arbitrary concurrent optical packets across multi-stage blocking networks, an optical RNS multiplier routes only \textit{a single optical pulse} per multiplication. Therefore, rather than routing light through deep multi-stage permutation fabrics, each active input waveguide requires only a dedicated binary switch tree to steer its optical pulse to the designated modular output port $j = (i \times B) \pmod{m}$.

\paragraph{\textbf{The 16th Waveguide ($WG_{16}$) and Fermat Arithmetic:}}
Rather than stopping at a 15-channel alphabet, JANUS introduces the \textbf{16th active waveguide ($WG_{16}$)} alongside the zero-gated reference channel ($WG_0$), establishing a 17-channel physical alphabet ($N_{\mathrm{alphabet}} = 17$):
\begin{equation}
\mathcal{A} = \{ WG_0, WG_1, WG_2, \dots, WG_{16} \} \longleftrightarrow \{ 0, 1, 2, \dots, 16 \}
\end{equation}
Because $17$ is a Fermat prime ($F_2 = 2^{2^1} + 1 = 17$), this topology unlocks major mathematical and physical properties:
\begin{enumerate}
\item \textbf{$100\%$ State Space Efficiency in $\mathbb{Z}_{17}$:} All 17 spatial channels represent valid residue classes. The full product matrix of $17 \times 17 = 289$ states is mapped with zero unused channels and zero register waste.
\item \textbf{Tree 16 Modular Negation Symmetry:} For any operand $W \in \mathbb{Z}_{17}$, multiplication by $16$ satisfies:
\begin{equation}
16 \times W \equiv -W \pmod{17}
\end{equation}
Tree 16 directly computes the additive inverse in $\mathbb{Z}_{17}$, eliminating modular subtraction logic.
\item \textbf{Single-Product Bound ($16 \times 16 = 256 < 257$):} The maximum element product never exceeds 256.
\item \textbf{Radix-16 $\mathbb{Z}_{257}$ Division-Free Reduction:} The next Fermat prime is $F_3 = 2^{2^2} + 1 = 257$. Because $16^2 = 256 \equiv -1 \pmod{257}$, any wide-word intermediate product $Y = 16 Y_H + Y_L$ reduces modulo $257$ via pure subtraction without integer division:
\begin{equation}
Y \pmod{257} = (Y_L - Y_H) \pmod{257}
\end{equation}
\end{enumerate}

\paragraph{\textbf{Binary Switch Tree Complexity and Stage Scaling:}}
Each of the 16 active waveguides ($WG_1 \dots WG_{16}$) connects to an independent 4-stage binary switch tree ($\log_2 16 = 4$ stages). Stage $s \in \{0, 1, 2, 3\}$ contains $2^s$ switches. The switch count per active input tree is:
\begin{equation}
N_{\mathrm{tree}} = \sum_{s=0}^{3} 2^s = 1 + 2 + 4 + 8 = 15~\mathrm{switches}
\end{equation}
Across all 16 active input trees, the multiplier core requires:
\begin{equation}
N_{\mathrm{switch,single}} = 16 \times 15 = \mathbf{240~\mathrm{switches}}
\end{equation}
The zero channel ($WG_0$) requires no dynamic switching elements (gated directly to $APD_0$). With an optical path depth of only $\mathbf{4~\mathrm{stages}}$, the optical fabric insertion loss is restricted to $1.61\,\mathrm{dB}$ and optical routing time-of-flight to just $1.33\,\mathrm{ps}$.

\begin{table*}[htbp]
\centering
\footnotesize
\setlength{\tabcolsep}{4pt}
\renewcommand{\arraystretch}{1.15}
\caption{COMPARATIVE EVALUATION OF INTEGRATED OPTICAL ROUTING ARCHITECTURES}
\label{tab:table5}
\begin{tabular*}{\textwidth}{@{\extracolsep{\fill}}lcccccc@{}}
\toprule
\textbf{Architecture / Topology} & \textbf{Switches/Mult} & \textbf{Stages ($S$)} & \textbf{Loss (dB)} & \textbf{Latency (ps)} & \textbf{Static Hold Power} & \textbf{State Efficiency} \\
\midrule
Full Crossbar ($256 \times 256$) & $65{,}536$ & $256$ & $>75.0$ & $>2500$ & Continuous Bias & $56.2\%$ \\
Dilated Spanke--Bene\v{s} ($256 \times 256$) & $130{,}560$ & $512$ & $>120.0$ & $>5000$ & Continuous Bias & $56.2\%$ \\
Path-Independent (PILOSS, $256 \times 256$) & $32{,}640$ & $256$ & $>50.0$ & $>2000$ & Continuous Bias & $56.2\%$ \\
Multi-Stage Bene\v{s} Mesh ($256 \times 256$) & $1{,}920$ & $15$ & $7.50$ & $750$ & Continuous Bias & $56.2\%$ \\
\textbf{Asymmetric 16-Tree Fermat Core (Ours)} & $\mathbf{240}$ & $\mathbf{4}$ & $\mathbf{1.61}$ & $\mathbf{1.33}$ & $\mathbf{0.00\,\mathrm{W}}$ (PCM) & $\mathbf{100\%}$ ($\mathbb{Z}_{17} / \mathbb{Z}_{257}$) \\
\bottomrule
\end{tabular*}
\end{table*}

\subsubsection{\textbf{Tree Crosstalk Suppression and Physical Isolation}}
In dense optical meshes, waveguide crossings and multi-path loops introduce coherent crosstalk. In the Asymmetric 16-Tree Fermat Core, optical propagation is strictly unidirectional and acyclic: light injected into any active waveguide traverses a dedicated binary expansion tree where inactive branches are decoupled via non-volatile $\mathrm{Sb_2S_3}$ crystalline switching ($\Delta n = 0.60$, extinction ratio $> 25\,\mathrm{dB}$). Measured signal-to-crosstalk ratio (SCR) reaches $18.96\,\mathrm{dB}$ worst-case and $52.16\,\mathrm{dB}$ mean, completely preventing false optical threshold detections at the APD receiver array.

\subsubsection{\textbf{Routing Loss Scaling Across Deployment Tiers}}
Each 16-Tree routing stage introduces a nominal insertion loss of $0.2627\,\mathrm{dB}$ (measured 3D FDTD). Across 4 stages:
\begin{equation}
L_{\mathrm{switch,4stg}} = 4 \times 0.2627\,\mathrm{dB} \approx 1.05\,\mathrm{dB}
\end{equation}
Including internal MMI splitters and crossing excess penalties, total fabric loss is $1.61\,\mathrm{dB}$.

\begin{table*}[htbp]
\centering
\small
\setlength{\tabcolsep}{6pt}
\renewcommand{\arraystretch}{1.15}
\caption{ROUTING LOSS AND SWITCH COMPLEXITY ACROSS JANUS DEPLOYMENT TIERS}
\label{tab:table6}
\begin{tabular*}{\textwidth}{@{\extracolsep{\fill}}lccc@{}}
\toprule
\textbf{Metric} & \textbf{Mini (32 Tiles)} & \textbf{Edge (32 Tiles)} & \textbf{Datacenter (32 Tiles)} \\
\midrule
Active Tiles & 32 & 32 & 32 \\
Intra-Tile Multipliers & $32 \times 32$ & $64 \times 64$ & $128 \times 128$ \\
Total Multipliers & $32{,}768$ & $131{,}072$ & $524{,}288$ \\
16-Tree Stages ($S$) & 4 & 4 & 4 \\
Total Switches (16-Tree Core) & $\mathbf{7{,}864{,}320}$ & $\mathbf{31{,}457{,}280}$ & $\mathbf{125{,}829{,}120}$ \\
Total Routing Loss & $\mathbf{1.61\,\mathrm{dB}}$ & $\mathbf{1.61\,\mathrm{dB}}$ & $\mathbf{1.61\,\mathrm{dB}}$ \\
Static Hold Power & $\mathbf{0.00\,\mathrm{W}}$ & $\mathbf{0.00\,\mathrm{W}}$ & $\mathbf{0.00\,\mathrm{W}}$ \\
Thermal Stability & Bounded Passive & Bounded Passive & Bounded Passive \\
\bottomrule
\end{tabular*}
\end{table*}

\subsubsection{\textbf{Architectural Properties of 16-Tree Fermat Routing}}
The Asymmetric 16-Tree Fermat Core achieves remarkable scaling characteristics:
\begin{enumerate}
\item \textbf{Monolithic Lithographic Feasibility:} Requires only 240 switches per multiplier, scaling from 3.93M in the 16-tile Model~1A prototype ($16,384$ multipliers) and 7.86M in the 32-tile Mini deployment ($32,768$ multipliers) up to 125.8M in Datacenter, well within single-reticle 3D manufacturing limits.
\item \textbf{High Optical Headroom:} Achieves an ultra-low fabric loss of $1.61\,\mathrm{dB}$, preserving a robust $+6.142\,\mathrm{dB}$ optical link margin without in-line amplification.
\item \textbf{Picosecond Core Latency:} Restricts optical time-of-flight across the switch fabric to $1.33\,\mathrm{ps}$, enabling sub-100~ps single-computation latency.
\item \textbf{Algebraic Completeness:} Maps $\mathbb{Z}_{17}$ with $100\%$ state space efficiency ($289/289$ states) and enables division-free Radix-16 $\mathbb{Z}_{257}$ modular reduction.
\end{enumerate}

\subsection{\textbf{High-Power Passive Fan-Out and Avalanche Photodetection}}
\label{sec:passive_fanout_apd}

After eliminating analog accumulation and static routing power through One-Hot spatial encoding and non-volatile PCM switching, the next dominant scaling constraint becomes optical propagation loss. Across multi-millimeter silicon photonic distribution trees and high-radix splitters, optical power attenuates significantly. Without rigorous link budget engineering and adequate initial launch power, the signal intensity can decay below reliable detection thresholds at the photodetectors.

As established in Section~\ref{sec:loss_accumulation}, in-line optical amplifiers (SOAs) accumulate spontaneous emission noise ($\sigma_{\mathrm{ASE}}^2 \propto G n_{\mathrm{sp}} B$) that destroys exact residue recoverability, while drawing unmanageable electrical power ($>30\,\mathrm{kW}$). JANUS structurally eliminates all intermediate optical gain elements, closing the optical link through high launch power passive fan-out and high-sensitivity avalanche photodetection.

\subsubsection{\textbf{Wavelength Selection: 1064 nm Yb-Fiber vs. 1550 nm Erbium}}
Surviving high passive fan-out losses without intermediate amplification requires multi-watt launch powers. At the conventional telecommunications wavelength of 1550 nm, Erbium-doped fiber amplifiers (EDFAs) exhibit a fundamental quantum defect efficiency limit dictated by the pump-to-signal photon energy ratio:

\begin{equation}
\eta_{QD,Er} = \frac{\lambda_{pump}}{\lambda_{signal}}
= \frac{976 nm}{1550 nm} \approx 62.9 \%
\end{equation}

Thus, generating 100 W of continuous-wave optical power at 1550 nm requires dissipating over 60 W of pure quantum-defect heat directly within the doped core, inducing severe thermal lensing, mode instability, and optical damage.

In contrast, operating at 1064 nm using Ytterbium-doped large-core fiber lasers pumped at 976 nm dramatically compresses the quantum defect:

\begin{equation}
\eta_{QD,Yb} = \frac{\lambda_{pump}}{\lambda_{signal}}
= \frac{976 nm}{1064 nm} \approx 91.7\% 
\end{equation}

yielding an intrinsic quantum defect of only 8.3\%. Combined with near-unity quantum yield and advanced thermal dissipation architectures, modern Yb-doped fiber lasers achieve wall-plug efficiencies (WPE) exceeding 75\% and deliver continuous-wave optical powers well above 100 W with negligible thermal degradation \cite{Jeong2004}.

Because both silicon nitride ($Si_3N_4$) passive routing waveguides and lithium tantalate ($LiTaO_3$) electro-optic Pockels micro-ring routers exhibit broad transparency and sub-0.1 dB/cm propagation loss at 1064 nm, shifting the carrier wavelength from 1550 nm to 1064 nm enables high optical launch power (up to 100 W / +50 dBm) while preserving strict single-mode optical confinement and eliminating the thermal ceilings of C-band sources.

\subsubsection{\textbf{Passive Splitting Budget and Link Margin}}

To distribute the master 1064 nm carrier across the computational mesh, JANUS utilizes a passive 1 : $2^{18}$ splitting tree formed from
18 cascaded 1 : 2 multimode interference (MMI) splitters (N = 262,144 output branches). The ideal passive splitting loss is:

\begin{equation}
L_{split,ideal} = 10 log_{10}(2^{18}) \approx 54.19 \text{ dB.}
\end{equation}

Incorporating excess optical path losses: under the 4-stage Asymmetric 16-Tree Fermat Core, optical routing attenuation is restricted to just $1.61\,\text{dB}$ ($4 \times 0.2627\,\text{dB} + \text{crossings}$). Combining this with MMI excess splitting loss ($\approx 5.40\,\text{dB}$) and propagation penalties ($1.50\,\text{dB}$), total excess path loss is bounded to $L_{\text{excess}} = 5.40\,\text{dB} + 1.61\,\text{dB} + 1.50\,\text{dB} = 8.51\,\text{dB}$. The end-to-end optical distribution loss is therefore:
\begin{equation}
L_{\text{total}} = L_{\text{split,ideal}} + L_{\text{excess}} = 54.19\,\text{dB} + 8.51\,\text{dB} = 62.70\,\text{dB}.
\end{equation}
With a master Yb-fiber laser delivering $P_{\text{laser}} = 100\,\text{W}$ (+50.0~dBm) optical launch power, the delivered signal power at each terminal detector branch reaches $P_{\text{det}} = -12.70\,\text{dBm}$ (or $-16.90\,\text{dBm}$ at reduced laser launch in the 16-tile Model~1A monolithic planar processor), providing an abundant power margin over the APD receiver sensitivity threshold.

\subsubsection{\textbf{Ge/Si $\mathrm{SAC^2M}$ Avalanche Photodetection}}

To achieve reliable binary One-Hot discrimination across delivered optical power levels from $-12.70\,\text{dBm}$ (100\,W carrier) down to $-16.90\,\text{dBm}$ (Model~1A sign-off corner) and $-18.59\,\text{dBm}$ (extreme degradation boundary) across a 100~GHz clock rate (within the device's inherent 105~GHz receiver bandwidth) without preamplification, JANUS replaces conventional PIN photodetectors with high-speed Separate-Absorption-Charge-Cliff-Multiplication ($\mathrm{SAC^2M}$) Ge/Si avalanche photodetectors \cite{Chen2025}. 

By confining impact ionisation to a thin silicon multiplication cliff layer separated from the germanium absorption region, these devices achieve internal avalanche gain $M = 7$ with an effective impact ionisation ratio $k \approx 0.06$. This yields a gain-bandwidth product of 441~GHz and an exceptionally low McIntyre excess noise factor:

\begin{equation}
F(M) = M\left[1-(1-k)
\left(\frac{M-1}{M}\right)^2\right]
\approx 2.0
\end{equation}

Paired with a receiverless 5 fF charge-integration front end and a StrongARM regenerative latch, the receiver incorporates an electro-optic Injection-Locked Oscillator (ILO) comb-referenced clock that reduces pulse timing jitter to 50 fs rms. The receiver sensitivity for a target bit error rate of $\le 10^{-36}$ is derived from the integrate-and-dump $Q$-factor ($Q \approx 12.72$, with $\sigma_{\mathrm{TIA}} = 1.8\,\mu\mathrm{A}$, $M=7$, and responsivity $R \approx 0.8\,\mathrm{A/W}$):
\begin{equation}
P_{\mathrm{sens}} = \frac{Q \cdot \sigma_{\mathrm{TIA}}}{M \cdot R} = \frac{9.38 \times 1.8 \times 10^{-6}}{7 \times 0.8} \approx 3.02\,\mu\mathrm{W} = -25.20\,\mathrm{dBm}.
\end{equation}
Adding a conservative $2.0\,\mathrm{dB}$ implementation margin for excess noise, coupling loss, and APD gain non-uniformity yields the practical binary detection sensitivity:
\begin{equation}
P_{\text{sens,practical}} \approx -23.20\,\text{dBm} \approx 4.79\,\mu\mathrm{W}.
\end{equation}

Consequently, under the master 100\,W Yb-fiber carrier ($P_{\mathrm{det}} = -12.70\,\text{dBm}$), the system establishes an abundant $+10.50\,\text{dB}$ link margin over practical sensitivity ($P_{\text{sens,practical}} = -23.20\,\text{dBm}$). In the compact 2.21\,W Model~1A processor, nominal delivered power ($P_{\mathrm{det}} = -14.63\,\text{dBm}$) provides $+8.41\,\text{dB}$ margin (and $+6.142\,\text{dB}$ under the conservative $-16.90\,\text{dBm}$ sign-off corner). Even under the extreme unpassivated degradation corner of $-18.59\,\text{dBm}$, the architecture preserves a robust $+4.61\,\text{dB}$ link margin (a $> 2.88\times$ power safety factor over the $10^{-18}$ BER threshold):

\begin{equation}
\begin{aligned}
\mathrm{Margin}_{\mathrm{100W}} &= (-12.70\,\mathrm{dBm}) - (-23.20\,\mathrm{dBm}) = +10.50\,\mathrm{dB}, \\
\mathrm{Margin}_{\mathrm{Model\,1A}} &= (-16.90\,\mathrm{dBm}) - (-23.04\,\mathrm{dBm}) = +6.142\,\mathrm{dB}, \\
\mathrm{Margin}_{\mathrm{corner}} &= (-18.59\,\mathrm{dBm}) - (-23.20\,\mathrm{dBm}) = +4.61\,\mathrm{dB}.
\end{aligned}
\end{equation}

\subsubsection{\textbf{Architectural Consequences}}

Replacing active in-line amplification with passive high-power fan-out and APD detection yields four critical architectural advantages:

\begin{itemize}
\item Elimination of Active Phase-Locking: Because no intermediate pump waves or four-wave mixing processes occur within the routing fabric, all requirements for active phase distribution and sub-micrometer path-length matching are eliminated.
\item Structural Absence of ASE Noise: Without optical gain elements in the datapath, spontaneous emission noise is completely absent ( ${\sigma^{2}}_{ASE}$ = 0 ). Link noise is strictly bounded by the deterministic thermal and shot noise floors of the receiver TIA and APD.
\item Reduced 3D Stacking Complexity: Removing active amplifier pump routing layers and associated waveguide tapers simplifies the interlayer optical coupling stack, reducing fabrication risk and removing vertical thermal barriers.
\item Deterministic Link Budget: System reliability becomes a direct function of master laser launch power and APD quantum efficiency, both of which operate well within demonstrated experimental envelopes.
\end{itemize}

\subsection{\textbf{Single-Wavelength Coherent Spatial Fanout}}

After eliminating analog precision collapse, routing thermodynamics, and amplifier noise accumulation, the next major scaling contradiction emerges from wavelength-domain multiplexing. Many photonic accelerators attempt to increase computational density through Dense Wavelength-Division Multiplexing (DWDM), where multiple optical carrier frequencies simultaneously propagate through shared waveguides \cite{Cotter1999},\cite{Miscuglio2020}.

While spectrally multiplexed systems improve theoretical bandwidth density, they become increasingly unstable at large optical powers due to nonlinear inter-channel interactions. JANUS therefore abandons multi-wavelength scaling entirely and instead adopts a single-wavelength coherent spatial fanout architecture.

As detailed in Section~\ref{sec:nonlinear_failure}, dense wavelength multiplexing (DWDM) introduces severe inter-channel four-wave mixing ($\omega_{\mathrm{FWM}} = \omega_1 + \omega_2 - \omega_3$) and cross-phase modulation that scale superlinearly with optical power density ($\sigma_{\mathrm{NL}}^2 \propto \gamma^2 P^3 L_{\mathrm{eff}}^2 / A_{\mathrm{eff}}$). JANUS eliminates inter-channel nonlinear spectral interactions by operating on exactly one optical carrier wavelength ($\lambda_0 = 1064\,\mathrm{nm}, N_\lambda = 1$), scaling computational bandwidth exclusively through spatial waveguides.

\subsubsection{\textbf{Single-Wavelength Reformulation}}

JANUS resolves this contradiction by eliminating multi-wavelength interaction entirely. Instead of scaling through spectral multiplexing, JANUS scales spatially. The architecture operates using exactly one optical carrier frequency:

\begin{equation}
N_{\lambda} = 1
\end{equation}

Thus, the optical field reduces to:
\begin{equation}
E(t) = Ae^{j\omega t}
\end{equation}

Without multiple carrier frequencies, the conditions required for dominant inter-channel Four-Wave Mixing are substantially suppressed:

\begin{equation}
\omega_1 = \omega_2 = \omega_3
\end{equation}

Therefore:
\begin{equation}
\omega_{\text{FWM}} =  \omega
\end{equation}

and no parasitic inter-channel spectral components emerge. 

Similarly, inter-channel Cross-Phase Modulation is substantially reduced because no independent wavelength channels exist to modulate one another. Thus, JANUS suppresses dominant nonlinear inter-channel interference architecturally rather than algorithmically. The resulting system obeys:

\begin{equation}
P_{ghost} \rightarrow 0
\end{equation}
under ideal operation.

% [Table 7 qualitative comparison integrated into text]


\subsubsection{\textbf{Active Single-Wave Injection}}

Because JANUS uses only a single optical carrier, the optical source can be coupled efficiently into the computation fabric. Rather than passive splitting into all candidate channels simultaneously, JANUS employs an active $1 \times 17$ electro-optic router at each multiplier input (or a 4-stage binary selector addressing the 16 active waveguides, with the zero channel handled directly). Each electro-optic router is fabricated in lithium tantalate ($\mathrm{LiTaO_3}$), exploiting the linear Pockels electro-optic effect rather than carrier injection or thermo-optic tuning. 

Pockels-effect switching stores energy capacitively and discharges it fully each cycle, yielding an effective dynamic energy of $E_{\mathrm{router,eff}} \approx 156.2\,\mathrm{aJ}$ per multiplier injection at $100\,\mathrm{GHz}$ (incorporating CMOS driver parasitics). For the Datacenter configuration ($524{,}288$ active input routers across $32$ tiles), the total router dynamic power is:

\begin{equation}
\begin{aligned}
P_{\mathrm{router}} &= N_{\mathrm{mult}} \times E_{\mathrm{router,eff}} \times f_{\mathrm{clk}} \\
&= 524{,}288 \times 156.2 \times 10^{-18} \times 100 \times 10^{9} \approx 8.19~\mathrm{W}
\end{aligned}
\end{equation}

consistent with Table~\ref{tab:table15}.

In spatial one-hot encoding, input operands are translated directly into spatial channel selection. By employing active electro-optic routing at each multiplier input, the optical carrier is injected exclusively into the designated waveguide, preserving optical link margin and avoiding the energy overhead associated with passive broadcast-and-attenuate schemes.

The signal laser (external to the chip) is a high-power 1064 nm continuous-wave (CW) Ytterbium-doped fiber laser source delivering Plaser = 100 W optical output power with a wall-plug efficiency (WPE) exceeding 75\% \cite{Jeong2004}, corresponding to an electrical power draw of $\approx$ 133 W (Table~\ref{tab:table15}). 

Because intermediate optical amplification layers are removed from the optical datapath, no dedicated per-tile optical pump lasers are required during inference.

\paragraph{\textbf{High-Power Passive Fan-Out Tree:}}

To distribute the master 1064 nm carrier across the computational mesh without intermediate optical amplification, JANUS employs a passive
1 : $2^{18}$ optical splitting tree composed of 18 cascaded 1 : 2 MMI splitter stages ( N = $2^{18}$ = 262,144 output ports). The ideal passive splitting loss across $2^{18}$ branches is:

\begin{equation}
L_{split,ideal} = 10 log_{10}(2^{18}) = 18 \times 3.0103 \approx 54.19 dB.
\end{equation}

Incorporating excess optical path losses: under the 4-stage Asymmetric 16-Tree Fermat Core, routing loss is limited to $1.61\,\text{dB}$, establishing a total excess loss of $L_{\text{excess}} = 8.51\,\text{dB}$ (including MMI excess loss of $5.40\,\text{dB}$ and propagation penalties of $1.50\,\text{dB}$) and total path attenuation of:
\begin{equation}
L_{\text{total}} = 54.19\,\text{dB} + 8.51\,\text{dB} = 62.70\,\text{dB}.
\end{equation}
At $P_{\text{laser}} = 100\,\text{W}$ (+50.0~dBm) launch power, each terminal receiver branch receives $P_{\text{branch}} = -12.70\,\text{dBm}$, guaranteeing a positive detection margin ($+6.142\,\mathrm{dB}$ to $+10.50\,\mathrm{dB}$) over the Ge/Si APD practical sensitivity limit ($P_{\text{sens,practical}} \approx -23.20\,\text{dBm}$).

\paragraph{\textbf{Fan-Out Architecture Reconciliation Note:}}

We note that the 1 : $2^{18}$ (262,144-port) passive splitting tree describes the primary optical carrier distribution from a single master laser, whereas the full Datacenter tile architecture specifies 32 tiles $\times$ 16,384 multipliers $\times$ 17 waveguides = $8{,}912{,}896$ total detector ports (524,288 active per cycle). These two descriptions address different points within the hierarchical optical distribution network; formal reconciliation mapping each of the 262,144 passive carrier outputs to the active $1 \times 17$ $\mathrm{LiTaO_3}$ electro-optic input routers across tiles remains open engineering work, as discussed in Section~\ref{sec:open_challenges}. Because active single-wave injection routes only a single One-Hot channel per multiplier, no power is wasted on dark channels. Importantly, all signal channels remain phase coherent because they originate from the same master carrier.

Thus:
\begin{equation}
\Delta \phi_{\text{source}} \approx 0
\end{equation}

across fanout branches.
This substantially simplifies global synchronization and phase management.

\begin{table}[htbp]
\centering
\small
\setlength{\tabcolsep}{7pt}
\caption{SINGLE-WAVELENGTH ACTIVE SIGNAL INJECTION SCALING ACROSS JANUS CONFIGURATIONS}
\label{tab:table8}
\begin{tabular}{lccc}
\toprule
Metric & Mini & Edge & Datacenter \\
\midrule
Active Injection Loss & $\approx$ 0 dB & $\approx$ 0 dB & $\approx$ 0 dB \\
Carrier Frequencies & 1 & 1 & 1 \\
Inter-Channel FWM & Suppressed & Suppressed & Suppressed \\
Inter-Channel XPM & Suppressed & Suppressed & Suppressed \\
Phase Coherence & Maintained & Maintained & Maintained \\
\bottomrule
\end{tabular}
\end{table}

\subsubsection{\textbf{Spatial Scaling Instead of Spectral Scaling}}

Conventional photonic systems attempt to increase throughput by increasing spectral density:

\begin{center}
\text{T $\propto$ N$_{\lambda}$}
\end{center}

However, JANUS instead scales throughput spatially:

\begin{center}
\text{T $\propto$ N$_{spatial}$}
\end{center}

This distinction is fundamental. Spectral scaling introduces nonlinear interaction between channels. Spatial scaling preserves channel independence because each computation occupies a physically distinct waveguide path. 

Therefore:

\begin{equation}
\frac{\partial C_i}{\partial C_j} \approx 0  \forall i \neq  j
\end{equation}

where: $C_i$ represents computational state of spatial channel i. This independence is critical for maintaining exact One-Hot arithmetic integrity.

\subsubsection{\textbf{Optical Entropy Isolation}}

An important consequence of single-wavelength operation is entropy localization. In multi-wavelength systems, nonlinear mixing causes entropy generated within one channel to propagate into neighboring channels. JANUS substantially reduces this propagation because all
channels share identical carrier frequency and remain spatially separated.

Thus:
\begin{equation}
\mathcal{E}_{i \rightarrow j} \rightarrow 0
\end{equation}

where: $\mathcal{E}_{i\rightarrow j}$ represents inter-channel entropy transfer.

This property is especially important for exact residue arithmetic, where even weak cross-channel corruption can invalidate CRT reconstruction.

\subsubsection{\textbf{Architectural Consequences}}

The transition from spectral multiplexing to single-wavelength spatial fanout fundamentally changes the scaling behavior of photonic computation.

\begin{enumerate}
\item Dominant nonlinear inter-channel interference mechanisms are substantially suppressed.
\item Coherent synchronization complexity collapses because all channels originate from a common optical source.
\item Throughput scaling becomes spatial rather than spectral.
\item Entropy propagation between channels is structurally suppressed.
\item The optical system remains within bounded non-linear operating conditions regardless of routing scale.
\end{enumerate}

Thus, JANUS does not attempt to mitigate nonlinear wavelength interference. It structurally suppresses the dominant spectral interaction
mechanism.

\subsection{\textbf{Spatial Residue Reconstruction and CRT Fault Visibility}}

After eliminating analog precision collapse, routing thermodynamics, amplifier noise accumulation, and spectral nonlinearities, the remaining challenge becomes large-integer precision recovery. Exact arithmetic systems require reconstruction of high-precision results from physically distributed low-precision optical computations. Conventional architectures attempt this reconstruction through wide binary arithmetic pipelines, which remain vulnerable to silent numerical corruption under transient faults.

JANUS instead performs precision scaling through spatial Residue Number System (RNS) decomposition combined with deterministic Chinese Remainder Theorem (CRT) reconstruction. This architecture produces a structurally important detection property:
\textbf{Residue perturbations produce discontinuous, out-of-range projections under CRT reconstruction rather than subtle silent corruptions.}

\subsubsection{\textbf{Spatial Residue Decomposition}}

Rather than executing wide-word arithmetic monolithically, JANUS decomposes computation into multiple independent residue domains.

For modulus set:

\begin{equation}
\mathcal{M} = \{m_1, m_2, \dots, m_k\}
\end{equation}

a large integer $X$ is represented as:
\begin{equation}
X \rightarrow (x_1, x_2, \dots, x_k)
\end{equation}

where:
\begin{equation}
x_i = X \mod m_i
\end{equation}

Each residue channel operates independently within a physically isolated optical tile.

The total representable dynamic range becomes:

\begin{equation}
M = \prod_{i=1}^k m_i \approx 2^{8k}
\end{equation}

Thus, large integer precision emerges from distributed modular composition rather than monolithic analog accumulation. Physically, each optical tile implements a modular residue channel across the 17 spatial waveguides ($WG_0$ zero bypass + $WG_1\dots WG_{16}$ active states) of the Asymmetric 16-Tree Fermat Core, supporting modular bases $m_i \le 257$ (e.g., $m_i \le 256$ in standard 8-bit residue decomposition).

Therefore:
\begin{equation}
x_i \in [0, m_i - 1] \subseteq [0, 255]
\end{equation}

within every local residue domain. When operating in isolated single-tile mode ($k=1$), a tile directly computes exact 4-bit $\times$ 4-bit multiplications ($P=4$, yielding an 8-bit product $Z \le 225 < 256$). To support wide operands ($P \ge 8$), wide integers are decomposed across $k = \lceil \frac{2P}{8} \rceil = \lceil \frac{P}{4} \rceil$ independent 8-bit residue tiles in parallel.

This bounded modular property prevents analog precision degradation inside individual optical pathways.

\textbf{Signed Integer Handling:} The RNS decomposition above operates on unsigned integers. To accommodate signed
operands, JANUS applies a bias offset of $2^{P -1}$ before residue encoding, where P is the native operand bit-width:

\begin{equation}
x' = x + 2^{P -1}
\end{equation}

This maps the signed range [$-2^{P -1}$, $2^{P -1}$ - 1] bijectively onto the unsigned residue domain [0, $2^{P - 1}$], allowing the same optical routing fabric to handle both signed and unsigned operands without hardware modification. The CMOS CRT backend subtracts the accumulated bias from the reconstructed result before returning the final signed output.

\subsubsection{\textbf{Chinese Remainder Reconstruction}}

After optical computation completes, the CMOS backend reconstructs the global result using the Chinese Remainder Theorem \cite{Mohan2016}\cite{Wang2000}\cite{Xia2000}.

For pairwise co-prime moduli:
\begin{equation}
gcd(m_i, m_j) = 1 \forall i \neq  j
\end{equation}

the unique reconstructed integer becomes:

\begin{equation}
X = \sum_{i=1}^k x_i M_i N_i \bmod M, \quad M_i = \frac{M}{m_i}, \quad N_i = M_i^{-1} \bmod m_i.
\end{equation}

This reconstruction occurs entirely within CMOS logic after optical routing completes. Importantly, optical hardware never performs high-precision analog accumulation.

Instead:
\textbf{Photonics $\rightarrow$ Bounded Modular Routing}

while:
\textbf{CMOS $\rightarrow$ Exact Integer Reconstruction}

Thus, exact arithmetic emerges through cooperation between bounded optical transport and deterministic digital verification.

\subsubsection{\textbf{Failure of Graceful Degradation}}

Conventional approximate architectures tolerate transient arithmetic errors through graceful degradation. Under such systems:

\begin{equation}
X \rightarrow X + \epsilon
\end{equation}

where: $\epsilon$ is a small perturbation.

This behavior is dangerous for exact computation because errors remain numerically plausible. Thermal noise, radiation events, timing faults, or analog drift may therefore silently propagate corrupted results into higher-level computation.

The resulting system exhibits 

\textit{Silent epistemic corruption}

where incorrect outputs remain computationally believable. Such behavior is unacceptable for deterministic exact arithmetic.

\subsubsection{\textbf{CRT Fault Visibility}}

JANUS substantially suppresses graceful degradation through CRT sensitivity \cite{Xia2000}, \cite{Barsi1973,Yang2002}.

Suppose residue channel $j$ experiences a transient perturbation ($x_j \to x_j + \delta$). The reconstructed result becomes $X' = (X + \delta M_j N_j) \bmod M$. Because $M_j = M / m_j \gg 1$, the deviation scales proportionally with the global dynamic range:
\begin{equation}
|X' - X| \propto M_j \gg |\delta|.
\end{equation}

Reconstruction error does not degrade linearly --- it amplifies combinatorially across the global integer space:

\textbf{Small Local Residue Fault $\rightarrow$ Discontinuous Global CRT Inconsistency}

The CMOS backend therefore detects highly anomalous arithmetic states with high probability. Instead of silently returning an incorrect value, the architecture invalidates the operation when CRT consistency checks fail.

% [Table 9 qualitative comparison integrated into text]


\subsubsection{\textbf{Error Detectability Bounds}}

The detectability condition emerges from CRT inconsistency.
For valid reconstruction:
\begin{equation}
X \mod m_i = x_i
\end{equation}

Any corrupted residue produces inconsistency:
\begin{equation}
X ' \mod m_j \neq  x_j 
\end{equation}

Thus, the CMOS backend performs deterministic consistency verification during reconstruction.

The probability of undetected residue corruption becomes:
$P_{\mathrm{undetected}} \to 0$
under bounded fault assumptions.

Importantly, fault visibility improves with increasing CRT dynamic range.
Therefore, larger systems become:
\textbf{more fault visible rather than less.}

This behavior fundamentally differs from conventional analog systems, where increasing scale typically reduces observability.

\subsubsection{\textbf{Spatial Fault Isolation}}

Because each residue tile operates independently, fault propagation remains spatially localized before reconstruction.

A thermal spike, fabrication defect, or transient radiation event affects only one residue channel:
\begin{equation}
\Delta x_i \neq  \Delta x_j (i \neq  j)
\end{equation}

Thus, no analog carry-chain propagation exists between tiles. This isolation substantially improves diagnosability and operational stability. Instead of diffuse analog corruption across the entire computational fabric, JANUS produces sharply localized detectable failure domains.

\subsubsection{\textbf{Precision Scaling Through Physical Parallelism}}

Conventional binary arithmetic increases precision temporally through carry propagation and wide-word accumulation. JANUS instead scales precision spatially. Additional precision requires additional residue channels:

\begin{equation}
P_{precision} \propto N_{tiles}
\end{equation}

Thus, computational precision scales through physical parallelism rather than analog signal fidelity. This distinction is fundamental. Conventional systems require exponentially increasing analog precision to scale arithmetic width. JANUS maintains bounded local precision while expanding global dynamic range through distributed modular composition.

\subsubsection{\textbf{Hybrid Partitioning and 3-Equation 64-Bit Decomposition}}

Scaling spatial RNS to wide 64-bit integer operands (INT64) presents a fundamental dynamic range challenge if executed via a monolithic optical basis. A native 64-bit multiplication ($2^{128}$ states) exceeds the dynamic range of a standard 16-tile optical cluster.

To make exact 64-bit RNS computation physically viable without precision collapse, JANUS enforces a strict architectural dichotomy known as Hybrid Partitioning:

\begin{quote}
\textbf{Optics exclusively performs $1 \times 1$ element-wise modular multiplication. All wide-word accumulation (addition) occurs strictly in digital CMOS after binary StrongARM latch detection.}
\end{quote}

Because optics is restricted to multiplication, each active waveguide conveys a single bounded modular product. High-speed StrongARM regenerative latches resolve binary waveguide states directly (zero ADCs), passing registered residues into 160-bit binary carry-save accumulators and the Dual-LUT Cross-Term Engine to sum wide GEMM arrays.

\paragraph{\textbf{Dual-LUT Cross-Term Engine \& 3-Equation Decomposition:}}
Instead of scaling optical tile counts, JANUS partitions 64-bit operands into 32-bit halves ($X_H, X_L$ and $Y_H, Y_L$) and computes three distinct equations:

\begin{enumerate}
\item \textbf{Equation 1 ($X_L \cdot Y_L$):} Computed fully optically in \textbf{Cluster 1} (8 tiles).
\item \textbf{Equation 2 ($X_H \cdot Y_H$):} Computed fully optically in \textbf{Cluster 2} (8 tiles).
\item \textbf{Equation 3 ($X_L Y_H + X_H Y_L$):} Evaluated via the \textbf{Dual-LUT Cross-Term Engine}---using localized CMOS SRAM Lookup Tables (LUTs) that decompose the 4-variable modular cross-term into parallel 16-bit lookups, achieving a $31,250\times$ memory reduction ($36\,\mathrm{GB}$ naive memory compressed to $1.152\,\mathrm{MB}$ raw / $1.5\,\mathrm{MB}$ physical on-chip SRAM).
\end{enumerate}

\paragraph{\textbf{Moduli Set Verification \& Formal Z3 Bounds:}}
The 8-modulus optical PRNS set allocated for Equation 1 and Equation 2 consists of $m_i \in \{255, 253, 251, 247, 241, 239, 233, 229\}$. Formal Z3 theorem proving verifies that this set yields a total dynamic range of:
\begin{equation}
M_8 \approx 5.68 \times 10^{19}
\end{equation}

The 32-bit halves ($X_H \cdot Y_H$) have a maximum product of $2^{62} \approx 4.61 \times 10^{18}$. 
Because $M_8$ strictly encapsulates the sub-products ($5.68 \times 10^{19} > 4.61 \times 10^{18}$), the system guarantees zero matrix overflow. The 9th modulus ($227$) is evaluated entirely within the CMOS Dual-LUT engine, reserving optical routing exclusively for direct sub-products.
\subsubsection{\textbf{Architectural Consequences}}

Spatial CRT reconstruction fundamentally changes the fault semantics of large-scale photonic arithmetic.

\begin{itemize}
\item  First, exact arithmetic emerges from bounded local residue computation rather than unstable global analog accumulation.
\item  Second, computational precision scales spatially rather than temporally.
\item  Third, local residue faults become globally detectable and unambiguously identifiable through CRT inconsistency amplification.
\item  Fourth, silent numerical corruption is structurally suppressed.
\item  Finally, the architecture structurally prevents the generation of numerically plausible incorrect outputs.
\end{itemize}

\begin{center}
\textbf{JIR-Predictive RRNS Redundancy and Throughput Trade-Off: } 
\end{center}

The baseline throughput figures assume all k tiles are allocated to compute. The JIR (JANUS Interface Representator) is the central tile-management controller responsible for packet routing, tile allocation, power monitoring, and thermal trend tracking. When the JIR observes a rising thermal trend on a tile --- rather than waiting for a fault to manifest --- it proactively reallocates one compute tile per RRNS parity channel ahead of the predicted fault threshold. This predictive engagement means RRNS is activated on thermal trajectory, not on residue inconsistency after the fact. For the Datacenter INT4 configuration (32 tiles, k = 1):

\begin{itemize}
\item  0 RRNS pairs: 32 compute engines $\rightarrow$ 44,564.5 TMAC/s
(baseline; all tiles nominal)
\item  1 RRNS pair: 31 compute engines $\rightarrow$ 43,171.8 TMAC/s
(-3.1\%; one tile thermally elevated)
\item  4 RRNS pairs: 28 compute engines $\rightarrow$ 38,993.9 TMAC/s
(-12.5\%; four tiles under thermal watch)
\end{itemize}

This trade-off is deterministic, controllable, and bounded. Full-throughput operation is the default; RRNS redundancy is engaged only when the JIR predicts elevated thermal risk from sensor trends, restoring compute tiles to normal allocation once temperature stabilizes within the safe operating band.

JANUS does not merely improve fault tolerance --- it substantially increases fault visibility by amplifying residue inconsistency through CRT reconstruction, converting small local faults into globally detectable arithmetic anomalies. The JIR further elevates system reliability by engaging RRNS protection on thermal prediction rather than reacting to errors after they occur.

\subsection{\textbf{3D Z-Axis Integration and Thermal Decoupling}}
\label{sec:zaxis_thermal}

After resolving analog precision instability, routing thermodynamics, optical attenuation, spectral nonlinearities, and arithmetic fault visibility, the final dominant scaling constraint becomes three-dimensional thermal interaction.

Large-scale photonic systems cannot sustain stable operation if CMOS logic, optical routing strata, memory structures, and detector arrays share a common planar thermal environment.

Even moderate lateral heat diffusion perturbs silicon refractive index, destabilizes optical phase alignment, and introduces global routing drift.

JANUS therefore adopts a vertically partitioned 3D Z-axis stack that thermally decouples the optical computation layer from the dominant heat sources of the system \cite{Skacel2021,Chang2017}.

As derived in Section~\ref{sec:thermal_drift}, lateral thermal diffusion in planar architectures entangles CMOS digital switching with temperature-sensitive photonic phase stability ($dn/dT \approx 1.86 \times 10^{-4}\,\mathrm{K}^{-1}$). JANUS structurally decouples these planes through a vertically partitioned 3D $Z$-axis stack:

\subsubsection{\textbf{Vertically Partitioned Z-Axis Stack \& Full Package Height}}

JANUS resolves planar thermal coupling through strict vertical thermal partitioning. The architecture is organized into a nested structure: a bare-die active stack enveloped by a macro-scale thermal and EMI package. The active bare-die layer stacks the following functional
planes:

\textbf{CMOS $\rightarrow$ SiO$_2$ $\rightarrow$ SiPh strata $\rightarrow$ Detectors}

Each layer performs a physically isolated functional role \cite{Atabaki2018}, \cite{Fenouillet2009}. The CMOS substrate (50 $\mu$m) executes digital reconstruction and verification. Above it resides the Silicon Dioxide ($SiO_2$) thermal buffer layer (100 $\mu$m). Each silicon photonic routing stratum occupies 30 $\mu$m; the number of strata varies by configuration (Mini: 3 strata, Edge: 4 strata, Datacenter: 5 strata). Detector arrays (2 $\times$ 5 $\mu$m = 10 $\mu$m interleaved two-layer block) reside above the photonic layer.

The resulting per-tier active stack heights are:
\begin{itemize}
\item  Mini (3 SiPh strata): 50 + 100 + 3$\times$ 30 + 10 = 250 $\mu$m
\item  Edge (4 SiPh strata): 50 + 100 + 4$\times$ 30 + 10 = 280 $\mu$m
\item  Datacenter (5 SiPh strata): 50 + 100 + 5 $\times$ 30 + 10 = 310 $\mu$m
\end{itemize}

\paragraph{\textbf{Macro-Package and Heat Spreader Height:}}

The active die does not represent the final physical system height. Because the architecture relies on directional upward heat extraction, the active die is mechanically bonded to a dual-stage macro-package:

\begin{itemize}
\item Heat Spreader 1 (HS1): A 30 $\mu$m dense copper-pillar Micro-Matrix layer sitting directly above the detector layer.
\item Spreader Gap: A 50 $\mu$m spacing containing the ultra-high-density (950,000 $mm^{-2}$) Cu--Cu pillar array and structural air/vacuum gaps.
\item Heat Spreader 2 (HS2): The final 250 $\mu$m external convective slim-lid heat sink.
\end{itemize}

The macro-package adds 30 + 50 + 250 = 330 $\mu$m to the active stack. 
Thus, the total packaged heights per tier are:

\begin{equation}
\begin{aligned}
T_{package,Mini} = 250 + 330 = 580 \mu\text{m} \approx 0.58 mm \\
T_{package,Edge} = 280 + 330 = 610 \mu\text{m} \approx 0.61 mm \\
T_{package,Datacenter} = 310 + 330 = 640 \mu\text{m} \approx 0.64 mm \\
\end{aligned}
\end{equation}

This sub-millimeter Z-axis form factor enables direct adjacent-die stacking with HBM memory (typical HBM stack height: 720 $\mu$m), which is compatible with the 0.64 mm Datacenter package height.

\paragraph{\textbf{Internal SiPh Sub-Layer Partitioning:}}

The SiPh routing layer is implemented as a multi-stratum vertical stack to satisfy reticle-limit constraints. The total switch fabric for the Datacenter configuration comprises $125,829,120$ PCM switches ($240$ switches per multiplier across $524,288$ multipliers). Assuming a compact PCM switch footprint of $1.25\,\mu\text{m}^{2}$, the net active switch area occupies approximately $157.3\,\text{mm}^2$ in Datacenter (and only $4.92\,\text{mm}^2$ in Mini).

By partitioning the fabric into vertical sub-layers, the effective footprint per stratum remains comfortably within standard lithographic reticle limits ($26 \times 33\,\text{mm} = 858\,\text{mm}^2$). Adjacent sub-layers are optically coupled via vertical interlayer grating couplers and thermally isolated by thin SiO$_2$ spacer layers ($\sim 2\,\mu\text{m}$ each). This vertical staging ensures the routing fabric remains within a manufacturable single-reticle window, avoiding multi-die stitching that would compromise the optical path-length uniformity required for exact One-Hot arithmetic.

The physical footprint per configuration (accounting for routing waveguides, MMI fan-out, APD receivers, pads, and edge-exclusion overhead) is derived as follows:

\begin{itemize}
\item Mini (1--2 strata, $3.93\text{--}7.86\,\text{M}$ switches across 16--32 tiles): $\approx 25\text{--}40\,\text{mm}^{2}$
\item Edge (2--3 strata, $\approx 15.73\,\text{M}$ switches): $\approx 60\text{--}80\,\text{mm}^{2}$
\item Datacenter (3--4 strata, $\approx 125.83\,\text{M}$ switches): $\approx 200\text{--}300\,\text{mm}^{2}$
\end{itemize}

These dimensions ensure that the entire photonic compute fabric occupies a single-reticle footprint, facilitating thermal and mechanical uniformity across the entire Z-axis stack. The SiO2 thermal buffer occupies:

\begin{equation}
T_{SiO_2} = 100 \mu\text{m}
\end{equation}

\subsubsection{\textbf{Thermal Isolation Through Vertical Heat Routing:}}

The key architectural principle is directional thermal extraction: CMOS-generated heat preferentially propagates downward toward heat spreaders ($q_{\mathrm{CMOS}} \downarrow$) while photonic programming and detector heat propagate upward ($q_{\mathrm{photonic}} \uparrow$). The intermediate $\mathrm{SiO_2}$ layer acts as a thermal barrier suppressing cross-plane heat leakage, ensuring $\partial T_{\mathrm{SiPh}} / \partial T_{\mathrm{CMOS}} \to 0$ and maintaining $T_{\mathrm{SiPh}} < T_{\mathrm{critical}}$ under steady-state operation.



\subsubsection{\textbf{Quantitative Thermal Resistance Network and $\Delta$ T Derivation}}

\paragraph{\textbf{Perimeter Thermal Shunt and EMI Faraday Shield:}}

JANUS implements directional heat extraction through a perimeter copper thermal shunt surrounding the 300.0~$\text{mm}^{2}$ silicon die. This low-thermal-resistance perimeter conduit connects the bottom CMOS substrate directly to the upper primary heat spreader (HS1) along the exterior periphery of the package.

This continuous perimeter structure serves a multifunctional triple purpose:
\begin{enumerate}
\item \textbf{Directional Heat Routing:} It establishes a low-impedance lateral thermal conduit that routes CMOS digital dissipation away from the active optical core directly into the primary heat sink, isolating the SiPh routing and $\mathrm{SiO_2}$ buffer layers from digital thermal flux.
\item \textbf{EMI Faraday Shielding:} It encases the internal stack to suppress high-frequency electromagnetic interference (EMI) at the 100~GHz modulation rate.
\item \textbf{Power Delivery Network:} The high-conductivity solid copper volume functions as the primary low-impedance ground return path and power supply rail for the CMOS substrate, minimising IR drop across the chip's power grid.
\end{enumerate}

This structural multi-use of the perimeter copper shunt eliminates the need for a separate dedicated power distribution layer within the Z-axis stack.

The internal Cu--Cu bond topologies complement this perimeter shunt:

\begin{itemize}
\item  \textbf{CMOS/SiO2/SiPh Interfaces:} Feature a sparse bond density of 10,000 $mm^{-2}$, acting as a vertical thermal bottleneck to force CMOS heat outward into the perimeter thermal shunt.
\item  \textbf{SiPh to HS1:} Features a dense 900,000 $mm^{-2}$ Cu-Cu array, with gaps structurally filled by $SiO_2$, drawing optical heat upward.
\item  \textbf{HS1 to HS2:} Connected via a 950,000 $mm^{-2}$ density interface bridging the 50 $\mu$m macro-package gap.
\end{itemize}

The density asymmetry is deliberate. The 10,000 $mm^{-2}$ sparse interface at the CMOS/ SiO$_2$ boundary presents a high vertical thermal resistance ( $R_{th,down}$ = 0.488 K/W), forcing CMOS-generated heat to route laterally into the perimeter thermal shunt rather than vertically through the optical stack. Simultaneously, the 900,000 $mm^{-2}$ dense interface at the SiPh/HS1 boundary presents a low vertical resistance ( $R_{th,up}$ = 0 .227 K /W), creating a preferential upward extraction path for photonic layer heat. The two density values thus implement directional thermal isolation by hydraulic analogy: a pinched pipe below and an open pipe above. Because $R_{th,up}$ $\ll$ $R_{th,down}$, the majority of SiPh heat is efficiently pulled upwards into the package spreaders.

\paragraph{\textbf{Global $\mathrm{SiO_2}$ Thermal Buffer Mass and Transient Dynamics:}}
Instead of evaluating thermal transients on an isolated per-tile basis, the $100\,\mu\mathrm{m}$ $\mathrm{SiO_2}$ layer forms a monolithic global buffer spanning the entire $300.0\,\mathrm{mm}^2$ die. The total thermal mass and heat capacity of this global buffer ($\rho_{\mathrm{SiO}_2} = 2{,}200\,\mathrm{kg/m^3}$, $c_{p,\mathrm{SiO}_2} \approx 703\,\mathrm{J/(kg\cdot K)}$) are:
\begin{equation}
m_{\mathrm{SiO}_2} = A_{\mathrm{chip}} h_{\mathrm{SiO}_2} \rho_{\mathrm{SiO}_2} \approx 6.6 \times 10^{-5}\,\mathrm{kg}, \quad C_{\mathrm{SiO}_2} \approx 0.0464\,\mathrm{J/K}.
\end{equation}
During one JIR activation cycle ($\tau_{\mathrm{JIR}} = 5\,\mu\mathrm{s}$), even under worst-case full-die transient dissipation ($P_{\mathrm{total}} = 48\,\mathrm{W}$, $Q_{\mathrm{gen}} = 240\,\mu\mathrm{J}$), the adiabatic global temperature rise is strictly bounded:
\begin{equation}
\Delta T_{\mathrm{cyc}} = \frac{Q_{\mathrm{gen}}}{C_{\mathrm{SiO}_2}} = \frac{240\,\mu\mathrm{J}}{0.0464\,\mathrm{J/K}} \approx 0.0052\,^\circ\mathrm{C}\text{ per JIR cycle}.
\end{equation}
Under sustained full-load operation, the upward thermal extraction path dominates ($R_{\mathrm{th,up}} = 0.227\,\mathrm{K/W}$ vs. $R_{\mathrm{th,down}} = 0.488\,\mathrm{K/W}$), restricting steady-state SiPh rise to $\Delta T_{\mathrm{steady}} = P_{\mathrm{tile}} R_{\mathrm{th,up}} \approx 0.34\,^\circ\mathrm{C}$. This maintains an extraordinary $17\times$ safety margin below the critical phase excursion threshold:
\begin{equation}
\Delta T_{\mathrm{crit}} = \frac{\lambda_0 \phi_{\mathrm{max}}}{2\pi L \frac{dn}{dT}} = \frac{1064\,\mathrm{nm} \times \pi}{2\pi (500\,\mu\mathrm{m}) (1.86 \times 10^{-4}\,\mathrm{K}^{-1})} \approx 5.7\,^\circ\mathrm{C}.
\end{equation}
Furthermore, vertical heat diffusion across the baseline $100\,\mu\mathrm{m}$ buffer takes:
\begin{equation}
\tau_{\mathrm{diff}} = \frac{h_{\mathrm{SiO}_2}^2}{\alpha_{\mathrm{SiO}_2}} = \frac{(100\,\mu\mathrm{m})^2}{9.05 \times 10^{-7}\,\mathrm{m^2/s}} \approx 11.05\,\mathrm{ms} \gg \tau_{\mathrm{JIR}} = 5\,\mu\mathrm{s}
\end{equation}
(expanding to $\tau_{\mathrm{diff}} = 69.06\,\mathrm{ms}$ under the $250\,\mu\mathrm{m}$ monolithic $\mathrm{SiO_2}$ buffer implemented in the Model~1A processor die). Consequently, any CMOS-originated thermal transient is completely attenuated by the $\mathrm{SiO_2}$ buffer before it can propagate into the SiPh routing stratum.

\begin{table*}[htbp]
\centering
\small
\setlength{\tabcolsep}{5pt}
\caption{UNIFIED THERMAL PARAMETER SCALING ACROSS JANUS CONFIGURATIONS}
\label{tab:table11}
\begin{tabular}{lcccc}
\toprule
Parameter & Symbol & Mini & Edge & Datacenter \\
\midrule
Total chip footprint & $A_{\mathrm{chip}}$ & $40.0~\mathrm{mm^2}$ & $100.0~\mathrm{mm^2}$ & $300.0~\mathrm{mm^2}$ \\
Global SiO$_2$ thickness & $h_{\mathrm{SiO_2}}$ & $100~\mu\mathrm{m}$ & $100~\mu\mathrm{m}$ & $100~\mu\mathrm{m}$ \\
Global SiO$_2$ heat capacity & $C_{\mathrm{SiO_2}}$ & $6.2~\mathrm{mJ/K}$ & $15.5~\mathrm{mJ/K}$ & $46.4~\mathrm{mJ/K}$ \\
JIR cycle duration & $\tau_{\mathrm{JIR}}$ & $5~\mu\mathrm{s}$ & $5~\mu\mathrm{s}$ & $5~\mu\mathrm{s}$ \\
Global worst-case heat/cycle & $Q_{\mathrm{gen}}$ & $33~\mu\mathrm{J}$ & $83~\mu\mathrm{J}$ & $250~\mu\mathrm{J}$ \\
Max. global $\Delta T$/cycle & $\Delta T_{\mathrm{cyc}}$ & $0.0053~^\circ\mathrm{C}$ & $0.0054~^\circ\mathrm{C}$ & $0.0054~^\circ\mathrm{C}$ \\
Upward thermal resistance & $R_{\mathrm{th,up}}$ & $1.38~\mathrm{K/W}$ & $0.55~\mathrm{K/W}$ & $0.18~\mathrm{K/W}$ \\
Steady-state SiPh rise & $\Delta T_{\mathrm{steady}}$ & $0.41~^\circ\mathrm{C}$ & $0.38~^\circ\mathrm{C}$ & $0.32~^\circ\mathrm{C}$ \\
Phase-drift critical $\Delta T$ & $\Delta T_{\mathrm{crit}}$ & $5.7~^\circ\mathrm{C}$ & $5.7~^\circ\mathrm{C}$ & $5.7~^\circ\mathrm{C}$ \\
SiO$_2$ diffusion time & $\tau_{\mathrm{diff}}$ & $11.05~\mathrm{ms}$ & $11.05~\mathrm{ms}$ & $11.05~\mathrm{ms}$ \\
HS pillar density (top) & $\rho_{\mathrm{Cu,top}}$ & $900~\mathrm{k/mm^2}$ & $900~\mathrm{k/mm^2}$ & $900~\mathrm{k/mm^2}$ \\
CMOS pillar density (bottom) & $\rho_{\mathrm{Cu,bot}}$ & $10~\mathrm{k/mm^2}$ & $10~\mathrm{k/mm^2}$ & $10~\mathrm{k/mm^2}$ \\
\bottomrule
\end{tabular}
\end{table*}

The photonic layer occupies a thermodynamic operating region where temperature variation remains bounded within optical stability limits.
The permissible phase drift condition is:
\begin{equation}
|\Delta \phi| < \phi_{\text{max}}
\end{equation}

Substituting the thermal phase relation:

\begin{equation}
| \frac {2 \pi L}{\lambda} \frac{dn}{dT} \Delta T |
< \phi_{\text{max}}
\end{equation}

Thus, bounded temperature directly implies bounded routing stability. Because the SiPh layer is thermally insulated from large CMOS thermal fluctuations, the optical network remains inside a stable operating region:
\begin{equation}
H(t) \approx constant 
\end{equation}

throughout sustained operation.

This thermodynamic partition maintains the optical routing substrate within a tightly bounded thermal stability window ($\Delta T \le 0.048~\mathrm{K}$), ensuring that phase determinism is preserved under sustained workloads without requiring active temperature trimming.

\subsubsection{\textbf{Detector Layer Decoupling}}

Detector arrays are positioned above the photonic routing layer to minimize optical-electrical interference \cite{Steckler2024,Chen2025}.  This vertical placement provides two major advantages :

\begin{enumerate}
\item  Optical absorption occurs after routing completion, preventing detector thermal dissipation from perturbing upstream waveguides.
\item  Detector electrical activity remains physically separated from the sensitive photonic routing mesh.
\end{enumerate}

Thus:
\begin{equation}
C_{\text{detector} \rightarrow \text{routing}} \rightarrow 0 
\end{equation}

where: C represents thermoelectrical coupling.

\paragraph{\textbf{Interleaved Two-Layer Detector Block Layout:}}

Each photonic tile implements a full N $\times$ N matrix--vector multiplication, producing $N^{2}$ independent one-hot multiplication
results per clock cycle. To read every result, a photodetector must be placed at the terminus of every output waveguide of every multiplier:

\begin{equation}
N_{\mathrm{det}} = N_{\mathrm{tiles}} \times N_{\mathrm{mult/tile}} \times N_{\mathrm{waveguides}} = 32 \times 16{,}384 \times 17 = \mathbf{8{,}912{,}896}
\end{equation}

Although 8,912,896 detectors are physically present across the 17 spatial waveguides ($WG_0 \dots WG_{16}$), the total active detector count per full 10~ps clock cycle is only $N_{\text{active}} = N_{\text{tiles}} \times N_{\text{dim}}^2 = 524{,}288$, because One-Hot encoding guarantees exactly one waveguide per multiplier carries light. However, to conserve optical master laser power and eliminate an entire layer of passive splitting, the 10~ps ($100\text{ GHz}$) clock cycle is time-multiplexed into two discrete 5~ps phases. At any instantaneous moment (i.e., within a single 5~ps phase), only half of the photodetectors (262,144) are actively illuminated. This two-phase illumination strategy allows the optical energy to be steered sequentially rather than split simultaneously, functionally doubling the delivered optical power per receiver. Over the full 10~ps cycle, all 524,288 detection events are resolved. The baseline spatial activity ratio $\alpha = 1/17$ reduces dynamic detector power proportionally, as quantified in the power budget (Table~\ref{tab:table15}). The 105~GHz Ge/Si $\mathrm{SAC^2M}$ APD photodetectors occupy approximately $1.5\,\mu\text{m}^2$ per device; $8.91 \times 10^{6}$ detectors require only $\approx 13.37\,\text{mm}^2$ of total detector area within the 300~$\text{mm}^{2}$ Datacenter die footprint.

To accommodate this high device density without metal wiring bottlenecks, JANUS organizes the photodetector array into an interleaved two-layer detector block. Each detector block consists of two vertically stacked Ge/Si $\mathrm{SAC^2M}$ APD sub-arrays, each handling alternating odd and even waveguide channels, addressed through separate metal interconnect layers. This interleaved arrangement halves the required wiring pitch per layer, suppresses inter-detector electrical crosstalk, and improves the spatial uniformity of the thermal dissipation profile directly below the primary heat spreader (HS1). The two-layer structure is fabricated within the same backend module as the SiPh routing stack and preserves detector margin consistency under sustained high-throughput operation.

\subsubsection{\textbf{Event-Driven Dynamic Latched Readout}}

Scaling to $8.91 \times 10^{6}$ physical detectors requires extreme mitigation of static power dissipation. Conventional continuous-time transimpedance amplifiers (TIAs) consume unacceptable milli-watt-level static power per channel \cite{Achtenberg2025}, which would fundamentally break the thermodynamic scaling limit of the chip. To resolve this, JANUS eschews continuous linear readout in favor of an \textit{event-driven dynamic latched comparator} architecture (e.g., clocked StrongARM latches) integrated intimately with the Ge/Si APDs \cite{Modak2025,Claveau2024}. 

This topology trades continuous analog linearity for ultra-low-power binary sensing \cite{Zheng2025,Severini2023}. Dynamic regenerative latches consume essentially zero static power; dynamic power is consumed exclusively during state transitions triggered by a clock edge and an incoming photon pulse. Because JANUS operates on a strictly sparse One-Hot arithmetic basis, only 1 in 17 waveguides is optically active per cycle ($\alpha = 1/17$). Hierarchical wake-up gating ensures that the vast majority of latches remain asleep, and only the optically triggered channels resolve. Combining highly scaled integration (parasitic input capacitance $C_{\text{int}} \le 3\text{ fF}$) \cite{Laube2024} with column-shared sample-and-hold (S/H) multiplexing \cite{Mulberry2018} amortizes the already low per-decision latch energy ($\approx 10\text{--}100\text{ fJ}$) down to an effective electrical cost of $\approx 100\text{ aJ}$ ($0.1\text{ fJ}$) per active detection. As formalized in Table~\ref{tab:table15}, this bounds the Datacenter detector dynamic power to a highly manageable $5.24\text{ W}$, avoiding the thermal scaling wall inherent to continuous TIA arrays.

\subsubsection{\textbf{Mechanical and Fabrication Implications}}

The layered stack additionally improves fabrication modularity. Each subsystem can be optimized independently:
\begin{itemize}
\item  CMOS for logic density,
\item  SiPh for low-loss routing,
\item  detectors for responsivity and bandwidth.
\end{itemize}

Thus, no single fabrication process must simultaneously optimize all electrical, optical, thermal, and nonlinear properties.

This modularity substantially improves manufacturability relative to monolithic heterogeneous photonic integration.

\subsubsection{\textbf{Thermal Scaling Across Configurations}}

Thermal stability is preserved across all JANUS deployment scales because optical routing density remains vertically isolated from the
dominant CMOS heat plane. Importantly, thermal behavior does not catastrophically degrade with routing scale because the dominant heat paths remain directionally partitioned.



\subsubsection{\textbf{Architectural Consequences}}

The 3D Z-axis stack fundamentally changes the thermodynamic behavior of large-scale photonic systems.

\begin{enumerate}
\item Optical routing becomes thermally isolated from CMOS activity.
\item Phase drift remains bounded within stable operating limits.
\item Calibration overhead collapses because the transfer matrix remains approximately stationary.
\item Passive routing strata, photodetector arrays, and CMOS digital logic each occupy vertically partitioned, thermally decoupled physical planes.
\item Finally, the architecture achieves sustained thermodynamic equilibrium under high-throughput operation.
\end{enumerate}

Thus, JANUS does not merely cool a photonic accelerator. It reorganizes the entire computational stack around directional thermal stability.

\section{SYSTEM-LEVEL SCALING ANALYSIS}
\label{sec:system_scaling}

\subsection{\textbf{System-Level Scaling and Physical Equilibrium}}

Conventional accelerator architectures evaluate scalability using theoretical peak arithmetic density alone. For large-scale photonic tensor processors, however, physically sustainable throughput is strictly governed by simultaneous stability across thermal, optical, informational, and circuit regimes:
\begin{equation}
P_{\mathrm{noise}} < P_{\mathrm{signal}} < P_{\mathrm{sat}}, \quad T_{\mathrm{junction}} < T_{\mathrm{critical}}, \quad |\Delta\phi_{\mathrm{drift}}| < \phi_{\mathrm{max}}.
\end{equation}
Rather than scaling through dense wavelength multiplexing ($T \propto N_\lambda$), which incurs superlinear four-wave mixing noise ($\sigma_{\mathrm{NL}}^2 \propto \gamma^2 P^3 L_{\mathrm{eff}}^2 / A_{\mathrm{eff}}$) and severe cross-phase modulation, JANUS scales throughput strictly through spatial parallelism ($T \propto N_{\mathrm{spatial}}$). Each computation traverses a physically independent waveguide path ($\partial C_i / \partial C_j \approx 0$ for $i \neq j$), structurally confining inter-channel optical entropy.

Crucially, because the Asymmetric 16-Tree Fermat Core maintains a strictly invariant routing depth of $S = \log_2(16) = 4$ binary stages across all tile matrix dimensions ($N = 32$ to $128$), optical fabric insertion loss is bounded to $1.61\,\mathrm{dB}$. This preserves a robust $+4.61\,\mathrm{dB}$ to $+10.50\,\mathrm{dB}$ binary detection margin over the Ge/Si SAC$^2$M APD receiver threshold ($-23.20\,\mathrm{dBm}$) across all deployment tiers without intermediate optical amplification. 

Simultaneously, directional $Z$-axis heat extraction ($R_{\mathrm{th,up}} = 0.227\,\mathrm{K/W}$ vs. $R_{\mathrm{th,down}} = 0.488\,\mathrm{K/W}$) paired with $18.5\,\mathrm{kHz}$ JIR dynamic modulus rotation decouples digital CMOS dissipation from the optical routing strata ($\partial T_{\mathrm{SiPh}} / \partial T_{\mathrm{CMOS}} \to 0$), maintaining the photonic transfer matrix stationary under continuous execution.

\subsection{\textbf{System Power Budget}}

The total system power of JANUS is decomposed by functional layer. Because PCM routing switches are non-volatile, their static hold power is approximately zero under infrequent-reconfiguration workloads. Table~\ref{tab:table15} provides the full breakdown for all three deployment configurations under sustained full-load operation.

\begin{table}[htbp]
\centering
\small
\setlength{\tabcolsep}{1pt}
\caption{JANUS FULL-SYSTEM POWER BUDGET BY COMPONENT LAYER. PCM ROUTING SWITCHES CONTRIBUTE 0 W STATIC HOLD POWER. ALL VALUES ARE PEAK SUSTAINED ESTIMATES AT 100 GHz OPERATION.}
\label{tab:table15}
\begin{tabular}{lccc}
\toprule
Component Layer & Mini & Edge & Datacenter \\
\midrule
1064 nm Yb Laser ($75\%$ WPE) & $\approx 6.67\text{ W}$ & $\approx 26.67\text{ W}$ & $\approx 133.33\text{ W}$ \\
LiTaO3 Pockels Ring Routers & 1.02 W & 4.10 W & 8.19 W \\
Ge/Si APD Detectors ($M = 7$) & 0.33 W & 1.31 W & 5.24 W$^\dagger$ \\
Optical Amplification Layer & 0 W & 0 W & 0 W \\
PCM Routing Switches (static) & 0.00 W & 0.00 W & 0.00 W \\
CMOS (Encoders / Adders / CRT) & 2.01 W & 8.50 W & 37.03 W \\
JIR Scheduler and Control Logic & 3.00 W & 3.00 W & 3.00 W \\
\bottomrule
\multicolumn{4}{l}{\scriptsize $^\dagger$Assumes highly amortized event-driven S/H readout ($\approx 100\text{ aJ}$ per active detection).} \\
\end{tabular}
\end{table}

For reference, a thermally tuned equivalent system using conventional thermo-optic MZI heaters at the typical 0.10 mW per switch would require:

\begin{equation}
P_{heaters,conv} = N_{switch} \times 0.10 mW
\end{equation}

For the Datacenter configuration with $N_{\mathrm{switch}} = 125,829,120$ PCM switches ($240$ switches per multiplier across $524,288$ multipliers):

\begin{equation}
P_{\mathrm{heaters,conv}} = 125,829,120 \times 0.10 \times 10^{-3} = 12.58\,\mathrm{kW}
\end{equation}

By replacing thermo-optic MZIs with non-volatile PCM routing, JANUS achieves a static routing power reduction from over $12.58\,\mathrm{kW}$ to identically zero ($0.00\,\mathrm{W}$) across the entire Datacenter tier, eliminating the multi-kilowatt static thermal burden inherent to continuously biased interferometric meshes.

\subsection{\textbf{Peak and Sustained Throughput Derivation and RNS Precision Scaling}}
\label{sec:operating_frequency}

\paragraph{\textbf{Wave-Pipelined Operating Frequency Derivation:}}

In the JANUS wave-pipelined architecture, functional stages along the optical-to-electrical datapath are physically decoupled in space, operating concurrently on successive in-flight optical pulses. Consequently, the minimum clock cycle period $T_{\text{cycle}}$ is not the cumulative sum of stage latencies, but is bounded by the bottleneck single-stage processing and recovery duration:

\begin{equation}
T_{\text{cycle}} = \max\left( t_{\text{mod}}, \, t_{\text{PD}}, \, t_{\text{wire}}, \, t_{\text{guard}} \right) \le 10.0\,\text{ps}
\end{equation}

Each concurrent stage satisfies the 10.0~ps timing boundary:
\begin{enumerate}
\item \textbf{Electro-Optic Injection Interval ($t_{\text{mod}} \le 10.0\,\text{ps}$):} 
The $\mathrm{LiTaO_3}$ Pockels-effect input router formats and injects optical pulses with full-width half-maximum $\tau_{\text{FWHM}} \approx 3\text{--}5\,\text{ps}$ into the active One-Hot waveguide at a 10.0~ps repetition interval.

\item \textbf{Photodetector Carrier Clearance ($t_{\text{PD}} \approx 1.52\,\text{ps}$):}
\begin{equation*}
t_{\text{PD}} = \frac{1}{2\pi \cdot f_{3\text{dB}}} = \frac{1}{2\pi \times 105\,\text{GHz}} \approx 1.52\,\text{ps}
\end{equation*}
where the 105~GHz receiver bandwidth corresponds to the high-speed Ge/Si $\mathrm{SAC^2M}$ avalanche photodetectors \cite{Chen2025}. Photogenerated carriers clear the multiplication region in $\approx 1.52\,\text{ps}$, ensuring complete reset well before the next pulse arrives in the 10.0~ps slot.

\item \textbf{Electrical Interconnect Delay ($t_{\text{wire}} \approx 1.33\,\text{ps}$):}
Local on-chip wire propagation over a $200\,\mu\text{m}$ effective metal path ($v \approx 1.5 \times 10^8\,\text{m/s}$) requires:
\begin{equation*}
t_{\text{wire}} = \frac{200 \times 10^{-6}}{1.5 \times 10^8} \approx 1.33\,\text{ps}
\end{equation*}

\item \textbf{Inter-Pulse Guard Margin ($t_{\text{guard}} \approx 3.5\,\text{ps}$):}
Provides temporal isolation between adjacent in-flight optical wavefronts to prevent inter-symbol interference (ISI) and group velocity dispersion across the passive routing fabric.
\end{enumerate}

Thus, the sustained system operating clock frequency is:
\begin{equation}
f_{\text{clk}} = \frac{1}{T_{\text{cycle}}} = \frac{1}{10.0\,\text{ps}} = 100\,\text{GHz}
\end{equation}

\paragraph{\textbf{Pipelining and End-to-End Latency Clarification:}}

The 10.0~ps cycle time represents the throughput interval (the time between successive optical pulses entering the routing fabric), not the end-to-end latency of a single computation. The cumulative single-computation end-to-end latency consists of:

\begin{enumerate}
\item \textbf{PCM 16-Tree Fermat Core optical propagation:} The 4 routing stages ($\log_2 16 = 4$) have an ultrafast optical propagation latency of only $t_{\text{opt}} \approx 1.33\,\text{ps}$ across the compact binary tree fabric (measured in full-wave simulation). Because the non-volatile PCM states remain static throughout an inference epoch, successive optical pulses stream through the fabric every 10.0~ps in a wave-pipelined fashion.

\item \textbf{CMOS CRT adder-tree accumulation:} The CMOS backend accumulates and reconstructs residues with an optimized 8-stage pipelined adder-tree latency of approximately $t_{\text{CRT}} \approx 80.0\,\text{ps}$ (at $10.0\,\text{ps}$ per register stage, reduced from $210\,\text{ps}$). This latency is hidden behind deeply staged CMOS registers that process reconstruction in parallel with optical transport.
\end{enumerate}

The cumulative single-computation end-to-end latency is approximately:
\begin{equation}
T_{\text{latency}} \approx t_{\text{opt}} + t_{\text{PD}} + t_{\text{wire}} + t_{\text{CRT}} \approx 1.33 + 1.52 + 1.33 + 80.0 \approx 84.18\,\text{ps} < 100\,\text{ps}
\end{equation}
from initial pulse injection to final reconstructed word, achieving an ultrafast end-to-end latency and collapsing the required wave-pipelined queue depth down to only $8\text{--}9$ in-flight computation stages at 100~GHz.

\paragraph{\textbf{Per-Tile Raw MAC Rate:}}
Each photonic tile implements an $N \times N$ matrix--vector multiplication in one optical clock cycle. The raw per-tile multiply-accumulate rate is:

\begin{equation}
\dot{M}_{\text{tile}} = N^2 \times f_{\text{clk}}
\end{equation}

For the three JANUS deployment configurations:
\begin{equation}
\dot{M}_{\text{Mini,tile}} = 102.4\,\text{TMAC/s}, \quad \dot{M}_{\text{Edge,tile}} = 409.6\,\text{TMAC/s}, \quad \dot{M}_{\text{DC,tile}} = 1,638.4\,\text{TMAC/s}.
\end{equation}

\paragraph{\textbf{Sustained Utilization Factor:}}

Not all tile cycles correspond to valid computation. The JIR scheduler maintains tiles in a warm-dominant regime, with controlled transitions to Hot state for active computation. The sustained utilization factor is derived from the thermal duty-cycle analysis of Section~\ref{sec:zaxis_thermal}:

$\eta = 0.85$

This factor accounts for:

\begin{itemize}
\item JIR state-transition overhead ( $\approx$ 4 $\mu$s per rotation cycle),
\item PCM reconfiguration cycles (sparse, $\ll$ 1\% overhead under sustained inference),
\item  CRT reconstruction pipeline flushing (contained within the $\tau$JIR = 5 $\mu$s cycle, introducing no per-operation stall).
\end{itemize}

\subparagraph{\textbf{Simultaneous Weight Loading and Zero-Overhead Reconfiguration}}

PCM reconfiguration introduces far less than 1\% overhead because the JANUS architecture reprograms switches in parallel rather than sequentially. The entire weight fabric---up to $125,829,120$ non-volatile PCM switches in the Datacenter configuration ($240$ switches per multiplier across $524,288$ multipliers)---can, in principle, be reprogrammed in a single 0.1~ms window using electro-thermal micro-heaters integrated directly within the phase-change switch cells. During the weight-loading window, the underlying CMOS driver layer asserts parallel electrical programming current pulses across the 32 tile driver matrices via the vertical Cu--Cu interconnect grid. This electro-thermal addressing scheme delivers the required 10--50~pJ programming energy (SET/RESET) per switch without requiring any auxiliary optical write lasers or complex optical pump distribution networks.

Furthermore, the 0.1~ms weight-loading window is not idle for the CMOS backend. The accumulated CRT reconstruction backlog from the previous compute epoch is processed during this interval. Because the CMOS pipeline is deeply staged, the reconstruction of the final batch of optical results from the previous epoch overlaps entirely with the electrical programming pulse cycle, yielding zero effective CMOS stall overhead attributable to reconfiguration. The net system utilization therefore remains at $\eta = 0.85$ rather than degrading with model-weight update frequency.

\subparagraph{\textbf{Reconfiguration Energy Budget}}

PCM state transitions require write energy $E_{\text{PCM}}$ per switch, which varies with device geometry, active volume, and the crystallisation (SET) versus amorphisation (RESET) direction. For plasmonic PCM switches of the type employed in JANUS, three design tiers bound the operating envelope across $N_{\text{sw}} = 125,829,120$ switches ($\approx 1.258 \times 10^8$ switches):

\begin{itemize}
\item \textbf{Tier 1 (highly optimised, $\approx 100\,\mathrm{fJ}$):} $E_{\text{reconfig}}^{T1} = 1.258 \times 10^8 \times 100\,\mathrm{fJ} \approx 12.58\,\mu\text{J}$ ($P_{\text{transient}}^{T1} \approx 0.126\,\text{W}$).
\item \textbf{Tier 2 (balanced neuromorphic, $\approx 1\,\mathrm{pJ}$):} $E_{\text{reconfig}}^{T2} = 1.258 \times 10^8 \times 1\,\mathrm{pJ} \approx 0.126\,\text{mJ}$ ($P_{\text{transient}}^{T2} \approx 1.26\,\text{W}$).
\item \textbf{Tier 3 (conservative upper bound, $\approx 10\,\mathrm{pJ}$):} $E_{\text{reconfig}}^{T3} = 1.258 \times 10^8 \times 10\,\mathrm{pJ} \approx 1.26\,\text{mJ}$ ($P_{\text{transient}}^{T3} \approx 12.6\,\text{W}$).
\end{itemize}

Even under the conservative Tier 3 scenario, the total transient electrical programming power dissipation across the entire chip reaches only $12.6\,\text{W}$ during the 0.1~ms write window. Distributing this transient electrical load across the 300~$\text{mm}^{2}$ Datacenter die through the 32 spatially partitioned CMOS tile driver arrays yields an exceptionally benign peak transient power density of $\approx 0.042~\text{W/mm}^{2}$, which is well within the thermal extraction capability of the perimeter thermal shunt and remains far below the $\sim 1~\text{W/mm}^{2}$ limit demonstrated in high-performance silicon flip-chip packages. After the write window closes, the PCM state is latched non-volatilely and the transient power returns to zero; no sustained thermal load is imposed on the optical stack. In the operationally dominant mode---incremental reconfiguration updating only the fraction $f_{\Delta}$ of changed weights between inference epochs---the energy scales as $E_{\Delta} = f_{\Delta} \times E_{\text{reconfig}}$. For typical sparse fine-tuning updates ($f_{\Delta} < 1\%$), Tier 3 reduces to $< 13\,\mu\text{J}$ per update, making reconfiguration energy negligible relative to inference throughput. Full simultaneous weight loading across all switches is therefore characterized as an ``in-principle'' upper bound; operational reprogramming exploits incremental weight sparsity via the partitioned CMOS driver arrays. 

\paragraph{\textbf{Peak Theoretical Throughput ($\eta = 1.0$)}}
The peak theoretical throughput of the architecture represents the raw hardware computational capacity when all $N_{\text{tiles}}$ and optical multipliers operate concurrently at $f_{\text{clk}} = 100\text{ GHz}$ under 100\% utilization ($\eta = 1.0$):

\begin{equation}
\dot{M}_{\text{peak}} = N_{\text{tiles}} \times \dot{M}_{\text{tile}} = N_{\text{tiles}} \times N_{\text{dim}}^2 \times f_{\text{clk}}
\end{equation}

For single-tile direct operation ($P=4$, $k=1$):
\begin{equation}
\dot{M}_{\text{peak,Mini}} = 3,276.8\,\text{TMAC/s}, \quad \dot{M}_{\text{peak,Edge}} = 13,107.2\,\text{TMAC/s}, \quad \dot{M}_{\text{peak,DC}} = 52,428.8\,\text{TMAC/s}.
\end{equation}

\paragraph{\textbf{Sustained Operational Throughput ($\eta = 0.85$)}}
Accounting for JIR scheduling, thermal rotation, and CRT reconstruction pipeline flushing ($\eta = 0.85$), the sustained operational throughput is:

\begin{equation}
\dot{M}_{\text{sust}} = \eta \times \dot{M}_{\text{peak}}
\end{equation}

Thus, at single-tile direct operation ($P=4$, $k=1$):
\begin{equation}
\dot{M}_{\text{sust,Mini}} \approx 2,785.3\,\text{TMAC/s}, \quad \dot{M}_{\text{sust,Edge}} \approx 11,141.1\,\text{TMAC/s}, \quad \dot{M}_{\text{sust,DC}} \approx 44,564.5\,\text{TMAC/s}.
\end{equation}

\paragraph{\textbf{Precision Scaling via Spatial RNS ($P \ge 8$, $k \ge 2$)}} 

Each optical tile implements a modular residue channel of width $\log_2(256) = 8\text{ bits}$ ($m_i \le 256$). When multiplying two wide integers of operand width $P$ bits, the resulting product spans $2P$ bits of dynamic range. To reconstruct this product without modular ambiguity, the computation is partitioned across $k$ independent 8-bit residue tiles:

\begin{equation}
k = \left\lceil\frac{2P}{8} \right\rceil = \left\lceil \frac{P}{4} \right\rceil
\end{equation}

optical tiles allocated in parallel.

Thus:
\[
\begin{aligned}
k_{\text{INT4 }(P=4)} &= 1, \quad
k_{\text{INT8 }(P=8)} = 2, \quad
k_{\text{INT16 }(P=16)} = 4,\\
k_{\text{INT32 }(P=32)} &= 8, \quad
k_{\text{INT64 }(P=64)} = 16
\end{aligned}
\]

The precision-scaled peak and sustained throughput scale inversely with $k$:
\begin{equation}
\dot{M}_{\text{peak}}(P) = \frac{\dot{M}_{\text{peak,INT4}}}{k}, \qquad \dot{M}_{\text{sust}}(P) = \frac{\dot{M}_{\text{sust,INT4}}}{k}
\end{equation}

This linear inverse scaling in $k$ represents the exact cost of achieving higher arithmetic precision through spatial modular composition. For 64-bit precision ($k=16$), the 16 allocated tiles operate under Hybrid Partitioning, splitting the operands into 32-bit halves. The tiles compute the direct sub-products ($X_L \cdot Y_L$ and $X_H \cdot Y_H$) fully optically in two 8-tile clusters, while the massive cross-term ($X_L Y_H + X_H Y_L$) is offloaded entirely to CMOS SRAM Memory Lookup Tables (LUTs). This resolves the full 64-bit matrix multiplication without modular overflow. No analog signal quality or dynamic range is sacrificed; the tradeoff is purely in physical tile allocation and CMOS integration.

\paragraph{\textbf{Energy Efficiency at 64-bit Precision}}

\begin{equation}
\eta_{EE} = \frac{\dot{M}_{\text{sust}}(64\text{-bit})}{P_{\text{total}}}
\end{equation}

For the three configurations:
\begin{equation}
\eta_{EE,\mathrm{Mini}} \approx 13.36\,\text{TMAC/s/W}, \quad \eta_{EE,\mathrm{Edge}} \approx 15.98\,\text{TMAC/s/W}, \quad \eta_{EE,\mathrm{DC}} \approx 14.91\,\text{TMAC/s/W}.
\end{equation}

\subsection{\textbf{Comprehensive Multi-Precision Performance Profiles}}
\label{sec:multi_precision_profiles}

Table~\ref{tab:table16} presents the complete peak theoretical and sustained operational throughput scaling, along with energy efficiency figures, across all three JANUS hardware deployment configurations (Mini, Edge, and Datacenter), spanning native INT4 direct mode ($P=4, k=1$) up to INT64 exact spatial RNS mode ($P=64, k=16$). Peak theoretical throughput corresponds to 100\% hardware utilization ($\eta = 1.0$), while sustained operational throughput incorporates JIR thermal rotation duty cycles and pipeline management ($\eta = 0.85$). Throughput and energy efficiency are quantified in $\text{TMAC/s}$ ($1\text{ MAC} = 2\text{ FLOPS}$) and $\text{TMAC/s/W}$, respectively.

\begin{table*}[htbp]
\centering
\small
\setlength{\tabcolsep}{4pt}
\caption{JANUS MULTI-PRECISION PERFORMANCE PROFILES ACROSS MINI ($32 \times 32$), EDGE ($64 \times 64$), AND DATACENTER ($128 \times 128$) CONFIGURATIONS (32 TILES EACH)}
\label{tab:table16}\label{tab:table17}\label{tab:table18}
\begin{tabular}{lc|ccc|ccc|ccc}
\toprule
& & \multicolumn{3}{c|}{\textbf{Mini (32 Tiles, $32 \times 32$)}} & \multicolumn{3}{c|}{\textbf{Edge (32 Tiles, $64 \times 64$)}} & \multicolumn{3}{c}{\textbf{Datacenter (32 Tiles, $128 \times 128$)}} \\
\textbf{Precision} & $k$ & \textbf{Peak} & \textbf{Sust.} & \textbf{EE (TMAC/s/W)} & \textbf{Peak} & \textbf{Sust.} & \textbf{EE (TMAC/s/W)} & \textbf{Peak} & \textbf{Sust.} & \textbf{EE (TMAC/s/W)} \\
\midrule
INT4 Direct ($P=4$) & 1 & 3,276.8 & 2,785.3 & 251.5--213.8$^\dagger$ & 13,107.2 & 11,141.1 & 300.8--255.7 & 52,428.8 & 44,564.5 & 280.7--238.6 \\
INT8 Exact ($P=8$) & 2 & 1,638.4 & 1,392.6 & 125.7--106.9 & 6,553.6 & 5,570.6 & 150.4--127.8 & 26,214.4 & 22,282.2 & 140.3--119.3 \\
INT16 Exact ($P=16$) & 4 & 819.2 & 696.3 & 62.9--53.4 & 3,276.8 & 2,785.3 & 75.2--63.9 & 13,107.2 & 11,141.1 & 70.2--59.7 \\
INT32 Exact ($P=32$) & 8 & 409.6 & 348.2 & 31.4--26.7 & 1,638.4 & 1,392.6 & 37.6--32.0 & 6,553.6 & 5,570.6 & 35.1--29.8 \\
INT64 Exact ($P=64$) & 16 & 204.8 & 174.1 & 15.7--13.4 & 819.2 & 696.3 & 18.8--16.0 & 3,276.8 & 2,785.3 & 17.5--14.9 \\
\bottomrule
\multicolumn{11}{p{0.96\textwidth}}{\scriptsize Peak and sustained throughput in TMAC/s ($1\,\mathrm{MAC}=2\,\mathrm{FLOPS}$). Sustained incorporates JIR duty cycle ($\eta = 0.85$). $^\dagger$Full-system wall-plug power ($13.03\,\mathrm{W}$). For comparison, Model~1A 16-tile planar delivers $1,392.6\,\mathrm{TMAC/s}$ (INT4, $225.71\,\mathrm{TMAC/s/W}$) and $87.0\,\mathrm{TMAC/s}$ (INT64, $14.10\,\mathrm{TMAC/s/W}$) at $6.17\,\mathrm{W}$ on-chip package power.} \\
\end{tabular}
\end{table*}
\subsection{\textbf{Asymptotic Architectural Implications}}

The scaling behavior of JANUS differs fundamentally from conventional analog photonic accelerators.
Conventional architectures exhibit:

\textbf{Higher Scale $\rightarrow$ Higher Instability}

because thermal drift, nonlinear interaction, calibration overhead, analog precision requirements, and fault observability all worsen with growth. JANUS instead enforces bounded local operating conditions while scaling globally through modular spatial composition.

\textbf{Higher Scale $\rightarrow$ Higher Parallelism}

without proportionally increasing local physical instability.

This bounded-local-state principle is the primary architectural reason why JANUS remains physically scalable under exact arithmetic workloads.

\section{DISCUSSION}
\label{sec:discussion}

\subsection{\textbf{Reframing Photonic Tensor Processing: Spatial Localization vs. Analog Accumulation}}

Most photonic accelerator research attempts to maximize compute density by directly emulating electronic analog MAC operations in the optical domain \cite{Sorger2018,Feldmann2021,Miscuglio2020,Zhang2025}. However, our analysis demonstrates that analog accumulation encounters an insurmountable physical barrier: unreduced accumulation of wide matrices requires an exponential signal-to-noise ratio ($\mathrm{SNR} \propto 2^N$, exceeding $138\,\mathrm{dB}$ for $128 \times 128$ INT8). No physical optical receiver or ADC can simultaneously sustain such dynamic range and low noise at multi-GHz frequencies.

JANUS resolves this fundamental scaling failure by abandoning analog amplitude recovery entirely. Instead of estimating continuous optical intensity (``How much light exists?''), the photodetector answers a binary spatial question (``Which waveguide contains light?''). This collapses the required detection threshold from $138.4\,\mathrm{dB}$ analog SNR down to a binary threshold ($>0\,\mathrm{dB}$ SNR, $P_{\mathrm{th}} \approx -23.20\,\mathrm{dBm}$), decoupling arithmetic precision from optical attenuation and enabling deterministic exact integer computation.

\subsection{\textbf{The Photonic-Electronic Division of Labor}}

Rather than attempting all-optical computation---which encounters severe physical limitations in non-linear state storage and control logic---JANUS enforces a disciplined hybrid division of labor:
\begin{equation}
\text{\textbf{Photonics}} \longrightarrow \text{Picosecond Spatial Transport}, \quad \text{\textbf{CMOS}} \longrightarrow \text{Deterministic Integer Verification}.
\end{equation}
Photonics excels at high-bandwidth, low-latency spatial routing ($1.33\,\mathrm{ps}$ time-of-flight per tile) with zero static hold power. Digital CMOS excels at deterministic verification, Garner mixed-radix CRT reconstruction, and fault monitoring. Through Hybrid Partitioning, optics evaluates bounded $1 \times 1$ modular sub-products ($X_L Y_L$ and $X_H Y_H$), while 160-bit carry-save accumulators and the Dual-LUT Cross-Term Engine resolve the cross-terms and wide GEMM summations without optical dynamic range explosion.

\subsection{\textbf{CRT Arithmetic Fault Visibility vs. Silent Data Corruption}}

A fundamental hazard in conventional approximate photonic accelerators is silent data corruption (SDC), where small analog drifts produce plausible but incorrect outputs ($X \to X + \epsilon$). Under CRT reconstruction, however, arithmetic faults exhibit extreme sensitivity: any single-residue perturbation ($x_j \to x_j + \delta$) amplifies across the global coprime product ($|X' - X| \propto M_j \gg |\delta|$), projecting corrupted states into discontiguous, out-of-range integer values. The CMOS backend instantly invalidates the operation, converting hidden numerical degradation into unambiguous, hardware-actionable faults.

\subsection{\textbf{Limitations and Open Research Challenges}}
\label{sec:open_challenges}

Although JANUS resolves the primary scaling contradictions of analog optical computing, several integration challenges must be addressed as the architecture advances toward physical fabrication:
\begin{enumerate}
\item \textbf{In-Waveguide $\mathrm{Sb_2S_3}$ Switching Contrast at 1064 nm:} While $\mathrm{Sb_2S_3}$ exhibits wide bandgap transparency ($E_g = 1.72\,\mathrm{eV}$) and low unannealed film losses, high-extinction refractive index contrast ($\Delta n \approx 0.60$) and cycle endurance within slot-waveguide directional couplers under multi-watt continuous optical flux remain unverified in silicon foundry test chips.
\item \textbf{High-Density Ge/Si $\mathrm{SAC^2M}$ APD Array Integration:} Discrete $105\,\mathrm{GHz}$ SAC$^2$M APD devices have been demonstrated in literature \cite{Chen2025}, but vertical 3D bonding into monolithic arrays ($278,528$ detectors) with receiverless $1:1$ StrongARM latch coupling across thermal spacers remains experimentally uncharacterized.
\item \textbf{High-Power Master Laser RIN Closure:} Laser relative intensity noise (RIN) of the high-power Yb-fiber source represents a primary SNR contributor in the absence of intermediate optical gain stages; formal RIN-to-BER analytical and experimental closure across multi-stage splitters remains open.
\item \textbf{Multi-Stratum 3D Photonic Stacking:} While CMOS-compatible 3D silicon photonic platforms spanning two and three strata have been experimentally reported \cite{Feldmann2019}, scaling to the 5-stratum configuration required by the Datacenter fabric (accommodating $125.8\text{M}$ switches across a $251.7\,\mathrm{mm}^2$ footprint) requires wafer-scale characterization of vertical optical via yield and cumulative interlayer insertion loss.
\item \textbf{Formal Multi-Fault and Adversarial Bounds:} While single-residue corruptions produce unambiguous CRT invalidations, formal closed-form bounds on simultaneous multi-tile correlated faults (e.g., shared thermal events across adjacent channels) remain an active theoretical inquiry.
\end{enumerate}

\subsection{\textbf{Scope Boundaries: What JANUS Does Not Claim}}
\label{sec:scope_boundaries}

Intellectual honesty requires delineating the explicit scope boundaries of this architecture:
\begin{itemize}
\item \textbf{Not Perfect Fault Immunity:} JANUS does not claim absolute immunity to arbitrary hardware faults. While single-residue faults and high-magnitude perturbations are mathematically flagged, simultaneous adversarial errors confined strictly within identical residue equivalence classes could theoretically evade single-cycle detection without multi-channel RRNS parity.
\item \textbf{Component Power Table Envelope:} The system power model in Table~\ref{tab:table15} provides an illustrative, first-principles architectural envelope rather than a formally closed empirical post-silicon measurement.
\item \textbf{Not a General-Purpose Processor:} JANUS targets dense matrix-matrix tensor acceleration. It does not replace general-purpose digital CPUs or control-flow logic.
\item \textbf{Not an Elimination of Physical Noise:} JANUS does not eliminate physical noise or thermal fluctuations; rather, it replaces exponential dependence on analog amplitude fidelity with robust binary spatial threshold detection.
\item \textbf{Forward-Looking Systems Integration:} While foundational device components ($\mathrm{Sb_2S_3}$ PCM switches, SAC$^2$M APDs, LiTaO$_3$ modulators) are demonstrated in literature, the complete multi-million switch 3D stacked monolithic system represents an architecturally motivated integration blueprint.
\end{itemize}

\section{CONCLUSION}

This work introduced JANUS, a constraint-bounded photonic architecture for exact large-scale arithmetic computation. Unlike conventional analog photonic accelerators, JANUS avoids high-precision optical accumulation entirely and instead reformulates computation around One-Hot spatial routing, residue decomposition, single-wavelength coherent transport, non-volatile PCM Asymmetric 16-Tree Fermat switching ($WG_{16}$ extension), high-power passive fan-out with Ge/Si avalanche photodetection, and CRT-based exact reconstruction.

The architecture was derived through iterative elimination of physically unstable alternatives.
\begin{enumerate}
\item  Analog accumulation was shown to violate Shannon-limited information recovery requirements at high precision.
\item  Thermo-optic routing was shown to become thermodynamically non-scalable under large routing densities.
\item  Conventional SOA amplification was shown to accumulate irrecoverable ASE entropy.
\item  DWDM scaling was shown to induce nonlinear spectral corruption.
\item  Graceful degradation arithmetic was shown to permit silent data corruption.
\end{enumerate}

JANUS resolves these contradictions through structural architectural reformulation: binary spatial localization replaces analog magnitude recovery; single-wavelength coherent fanout removes inter-channel spectral mixing; non-volatile PCM routing eliminates static hold power; high-power 1064\,nm passive fan-out with Ge/Si avalanche photodetection survives routing loss without intermediate optical amplification; CRT reconstruction converts local arithmetic faults into globally detectable invalid states; and $Z$-axis thermal partitioning isolates the optical compute layer from dominant CMOS thermal activity. Collectively, these mechanisms ensure that the architecture remains consistent with Shannon information limits, thermodynamic stability, nonlinear optical bounds, bounded noise propagation, and deterministic fault visibility.

Importantly, JANUS does not claim to transcend physical law. Instead, it reorganizes computation so that all dominant physical constraints remain bounded simultaneously. The resulting system demonstrates that scalable photonic exact arithmetic is achievable not through increasingly aggressive analog precision, but through spatial modular reformulation aligned with the underlying laws of optics, thermodynamics, and information theory.

\appendices

\section{NOTATION AND SYSTEM PARAMETERS}
\label{app:notation}

\begin{table}[htbp]
\centering
\caption{PRIMARY NOTATION USED THROUGHOUT THIS WORK}
\label{tab:table22}
\scriptsize
\setlength{\tabcolsep}{2pt}
\begin{tabular}{ll|ll}
\toprule
Symbol & Definition & Symbol & Definition \\
\midrule
$N$ & Matrix/mesh dimension & $S$ & Routing stages ($S=4$ in 16-Tree) \\
$P_{\text{laser}}$ & Master laser optical power & $\phi$ & Optical phase \\
$L$ & Propagation/interaction length & $\alpha$ & Attenuation coeff. / diffusivity \\
$M$ & CRT product range: $M = \prod m_i$ & $T$ & Temperature or throughput \\
$M_j$ & CRT partial product: $M/m_j$ & $n$ & Refractive index \\
$m_i$ & $i$-th RNS modulus & $\chi^{(2)}$ & 2nd-order nonlinear susceptibility \\
$x_i$ & Residue of $X \bmod m_i$ & $\omega$ & Angular optical frequency \\
$N_i$ & Garner coefficient: $M_i^{-1} \bmod m_i$ & $\gamma$ & Nonlinear optical coefficient \\
$G, F$ & Avalanche gain, excess noise & $Q, \mathrm{BER}$ & Signal-to-noise, bit error rate \\
\bottomrule
\end{tabular}
\end{table}

\providecommand{\BIBdecl}{\relax}
\renewcommand{\BIBdecl}{\scriptsize\setlength{\itemsep}{-0.5pt}}
\bibliographystyle{IEEEtran}
\bibliography{references}

\end{document}
